\documentclass[12pt,a4paper,fontset=windows]{ctexart}
\usepackage[left=25mm,right=23mm,top=24mm,bottom=23mm,headheight=17pt]{geometry}
\usepackage[table]{xcolor}
\usepackage{graphicx,booktabs,longtable,tabularx,array,makecell}
\usepackage{enumitem,setspace,fancyhdr,titlesec,caption,float}
\usepackage{hyperref,bookmark,xurl,needspace,microtype}
\usepackage{amsmath,amssymb}
\usepackage{placeins}
\usepackage{pdfpages}
\emergencystretch=1.5em
\renewcommand{\topfraction}{.91}
\renewcommand{\textfraction}{.06}
\renewcommand{\floatpagefraction}{.55}
\setlength{\textfloatsep}{11pt plus 2pt minus 2pt}
\setlength{\floatsep}{9pt plus 2pt minus 2pt}
\setlength{\tabcolsep}{4pt}

\definecolor{Ink}{HTML}{20242A}
\definecolor{Muted}{HTML}{53606A}
\definecolor{Rule}{HTML}{D4DADD}
\definecolor{Pale}{HTML}{F3F5F6}
\hypersetup{colorlinks=true,linkcolor=Ink,urlcolor=Muted,citecolor=Ink,hypertexnames=false,
  pdfauthor={蒋宗达},pdftitle={安德鲁牙科连锁运营管理系统分析与设计}}
\setmainfont{Times New Roman}
\setsansfont{Segoe UI}
\setmonofont{Consolas}
\setCJKmainfont[AutoFakeBold=2]{SimSun}
\setCJKsansfont{Microsoft YaHei}
\setCJKmonofont{Microsoft YaHei}
\setstretch{1.22}
\setlength{\parindent}{2em}
\setlength{\parskip}{.12em}
\setlist{leftmargin=2.1em,itemsep=.12em,topsep=.25em}
\setcounter{tocdepth}{1}
\pagestyle{fancy}
\fancyhf{}
\fancyhead[L]{\small\sffamily\color{Muted}安德鲁牙科连锁运营管理系统}
\fancyhead[R]{\small\sffamily\color{Muted}分析与设计报告}
\fancyfoot[L]{\scriptsize\sffamily\color{Muted}蒋宗达\quad 810224025}
\fancyfoot[C]{\small\sffamily\color{Muted}\thepage}
\renewcommand{\headrulewidth}{.3pt}
\titleformat{\section}{\sffamily\bfseries\color{Ink}\fontsize{17}{22}\selectfont}{\thesection}{.7em}{}
\titleformat{\subsection}{\sffamily\bfseries\color{Ink}\fontsize{13}{18}\selectfont}{\thesubsection}{.65em}{}
\titleformat{\subsubsection}{\sffamily\bfseries\color{Ink}\normalsize}{\thesubsubsection}{.6em}{}
\titlespacing*{\section}{0pt}{17pt plus 3pt minus 2pt}{8pt}
\titlespacing*{\subsection}{0pt}{11pt plus 2pt minus 1pt}{4pt}
\captionsetup{font=small,labelfont={bf,color=Ink},labelsep=quad,skip=5pt}
\newcolumntype{L}[1]{>{\raggedright\arraybackslash}p{#1}}
\newcolumntype{C}[1]{>{\centering\arraybackslash}p{#1}}
\graphicspath{{../assets/uml/}{../assets/erd/}{../assets/erd/domains/}}
\newcommand{\figfull}[3]{\begin{figure}[H]\centering
  \includegraphics[width=\textwidth,height=.92\textheight,keepaspectratio]{#1.png}
  \caption{#2}\label{#3}\end{figure}}
\newcommand{\figshort}[3]{\begin{figure}[H]\centering
  \includegraphics[width=\textwidth,height=.78\textheight,keepaspectratio]{#1.png}
  \caption{#2}\label{#3}\end{figure}}
\newcommand{\figuc}[3]{\begin{figure}[H]\centering
  \includegraphics[width=.9\textwidth,height=.72\textheight,keepaspectratio]{#1.png}
  \caption{#2}\label{#3}\end{figure}}
\newcommand{\figdb}[3]{\begin{figure}[H]\centering
  \includegraphics[width=\textwidth,height=.77\textheight,keepaspectratio]{#1.png}
  \caption{#2}\label{#3}\end{figure}}
\newcommand{\figvector}[3]{\begin{figure}[H]\centering
  \includegraphics[width=\textwidth,height=.78\textheight,keepaspectratio]{#1.png}
  \caption{#2}\label{#3}\end{figure}}
\newcommand{\figclassvector}[3]{\begin{figure}[H]\centering
  \includegraphics[width=\textwidth,height=.78\textheight,keepaspectratio]{#1_vector.pdf}
  \caption{#2}\label{#3}\end{figure}}
\newcommand{\figvectorlong}[3]{\begin{figure}[H]\centering
  \includegraphics[width=\textwidth,height=.88\textheight,keepaspectratio]{#1.png}
  \caption{#2}\label{#3}\end{figure}}
\newcommand{\figstate}[3]{\begin{figure}[H]\centering
  \includegraphics[width=\textwidth,height=.45\textheight,keepaspectratio]{#1.png}
  \caption{#2}\label{#3}\end{figure}}
\newcommand{\figstatevector}[3]{\begin{figure}[H]\centering
  \includegraphics[width=\textwidth,height=.57\textheight,keepaspectratio]{#1.pdf}
  \caption{#2}\label{#3}\end{figure}}
\newcommand{\figoverview}[3]{\clearpage\includepdf[pages=1,fitpaper=true,pagecommand={\thispagestyle{empty}\refstepcounter{figure}\label{#3}}]{../assets/uml/overview_pages/#1.pdf}\clearpage}
\newcommand{\figanalysisplate}[3]{\clearpage\includepdf[pages=1,fitpaper=true,pagecommand={\thispagestyle{empty}\refstepcounter{figure}\label{#3}}]{../assets/uml/analysis_pages/#1.pdf}\clearpage}
\newcommand{\figerdplate}[3]{\includepdf[pages=1,fitpaper=true,pagecommand={\thispagestyle{empty}\refstepcounter{figure}\label{#3}}]{../assets/erd/pages/#1.pdf}}
\newcommand{\figcompact}[3]{\begin{figure}[H]\centering
  \includegraphics[width=.8\textwidth,height=.78\textheight,keepaspectratio]{#1.png}
  \caption{#2}\label{#3}\end{figure}}
\newcommand{\figdbportrait}[3]{\clearpage\begin{figure}[H]\centering
  \includegraphics[width=.98\textwidth,height=.86\textheight,keepaspectratio]{#1.png}
  \caption{#2}\label{#3}\end{figure}\clearpage}
\newcommand{\spec}[6]{\Needspace{10\baselineskip}\subsection{#1\quad #2}
  \noindent\textbf{执行者与前置条件}\quad #3\par
  \noindent\textbf{成功结果}\quad #4\par
  \noindent\textbf{基本路径}\quad #5\par
  \noindent\textbf{异常与规则}\quad #6\par}

\begin{document}
\begin{titlepage}
\centering
\vspace*{24mm}
{\sffamily\large 系统分析与设计课程设计\par}
\vspace{28mm}
{\sffamily\fontsize{26}{34}\selectfont\bfseries 安德鲁牙科连锁运营管理系统\par}
\vspace{7mm}
{\sffamily\LARGE 分析与设计报告\par}
\vspace{8mm}
{\large Andrew's Dental Chain\par}
\vfill
\renewcommand{\arraystretch}{1.55}
\begin{tabular}{@{}L{27mm}L{77mm}@{}}
学生姓名 & 蒋宗达\\
学号 & 810224025\\
课程名称 & 系统分析与设计\\
选题案例 & 安德鲁牙科连锁店：数据建模\\
日期 & 2026 年 10 月\\
\end{tabular}
\vspace*{19mm}
\end{titlepage}
\pagenumbering{Roman}
\section*{摘要}
\addcontentsline{toc}{section}{摘要}
安德鲁牙科连锁在广东经营 15 家诊所，计划在三年内扩大至约 30 家。随着服务网络扩展，分散的患者、预约与治疗记录增加了跨店复诊的信息核对成本，也使总部难以统一分析预约履约、专业能力配置和季度经营结果。本文以预约、到诊、诊疗、结算的连续过程为主线，建立业务流程、系统用例与分析模型，并将关键业务约束落实到服务职责和关系数据结构。

系统以稳定的全网患者编号贯通各店记录。预约提交时复核牙医任职、项目资质、服务时段及已分配牙椅；临床环节分别保存患者诉求、治疗计划和实际治疗，结算明细据实施项目形成。数据层包含 61 张关系表、730 个字段和 119 条现有外键，分为组织授权、身份通知、患者主档、人员目录、预约资源、临床就诊、结算资金、期间核算和扩张规划九个业务簇。季度核算保存收入、用品成本与运营费用的归属依据及核定快照。现有程序已覆盖主要登记、预约、临床与结算操作；候补邀约、牙医级费用分摊及正式短信、微信身份通道仍有待完善或联调。

\noindent\textbf{关键词：}牙科连锁；业务用例；预约管理；临床记录；诊疗结算；季度核算
\clearpage
\begingroup\small\setstretch{1.03}\tableofcontents\endgroup
\clearpage\pagenumbering{arabic}

\section{项目背景与愿景}
\subsection{问题来源}
案例中的 ADC 由 Andrew Chen 管理，正从现有 15 家诊所向更大规模发展。扩张带来四项管理问题：如何比较门店间的诊疗标准，识别取消预约的影响，跨店配置牙医的专业能力，以及核算牙医的季度贡献\cite{case}。回答这些问题的基础，是为患者、牙医、预约、就诊与结算建立一致的标识和记录口径，使汇总结果能够回溯到具体业务事实。

本文以课程案例界定业务范围。案例未说明门店实际使用的登记工具，现状流程中的本地档案、台账及人工移交据此作为待调研核实的分析假设。病史更新、治疗记录与风险复核的设计参考牙科执业资料\cite{ada-history,ada-records}；取消时限、知情同意和费用分摊等机构制度，须在正式实施前确认。

\subsection{愿景与目标}
ADC Nexus 面向前台、牙医、收费员、诊所经理、总部运营与财务人员，并为患者提供本人资料查询入口。系统以同一次就诊为线索连接预约、治疗和费用记录，同时保持岗位责任：前台核实身份，牙医作出治疗判断并签署记录，收费员核实收款，财务确认核算口径。软件承担检索、校验、记录和汇总职责，为专业判断提供可靠依据。

\begin{table}[!htbp]\centering\small
\caption{目标与评价办法}\label{tab:goals}
\begin{tabular}{@{}L{17mm}L{54mm}L{73mm}@{}}\toprule
目标 & 评价对象 & 拟采用的检查方法\\\midrule
G1 & 患者资料连续 & 授权跨店查询可见既往就诊；疑似重复记录有人工作出结论\\
G2 & 预约可执行 & 提交时同时检查时间冲突、牙医任职、资质与牙椅\\
G3 & 治疗与结算一致 & 结算行能追溯实际治疗；同一就诊不重复出单\\
G4 & 季度结果可复核 & 收入、成本及费用有归属依据，核定结果保留版本\\
\bottomrule\end{tabular}\end{table}

项目范围包括患者资料、预约与到诊、临床记录、逐项结算、人员资质、季度核算、经营查询及身份权限。保险理赔、法定税务开票、自动诊断和器械灭菌追踪不纳入本次设计。指标的改善幅度须在门店试点建立基线后测量，不能由合成数据推断。

报告依照课程课件的业务建模、需求、分析与设计工作流组织\cite{course}。业务流程描述岗位协作，系统用例规定用户与系统的交互，分析模型说明对象职责，设计部分再落实到服务边界、关系表和岗位界面。

\section{A 业务建模}
\subsection{业务用例}
业务用例以 ADC 组织为责任边界，患者是取得服务的外部执行者；前台、牙医、收费员及总部岗位的内部协作由业务序列图展开。图~\ref{fig:business-usecase} 从患者获得的结果出发，将服务归纳为诊疗、预约、结算与连续照护四项。
\figshort{01_业务用例图}{ADC 业务用例图}{fig:business-usecase}
图中的实线箭头从患者指向其参与的业务用例，仅标明参与方向，不表示四项服务的先后顺序；办理次序由后续业务序列图说明。

\begin{table}[!htbp]\centering\small
\caption{业务用例与完成结果}\label{tab:business-cases}
\begin{tabular}{@{}L{17mm}L{46mm}L{81mm}@{}}\toprule
编号 & 患者目标 & 业务完成结果\\\midrule
BUC-01 & 获得牙科诊疗服务 & 经检查和必要说明后接受相应治疗\\
BUC-02 & 安排与变更预约 & 确认诊所、牙医和时间；条件变化时取得变更或取消结果\\
BUC-03 & 完成费用结算 & 获得实际治疗的逐项结算单和收款凭据\\
BUC-04 & 获得连续照护 & 跨次、跨店就诊时可使用既往资料及复诊安排\\
\bottomrule\end{tabular}\end{table}

\subsection{现状业务流程}
现状模型包括单次就诊的预约、到诊、诊疗、结算四个阶段，以及独立发生的季度经营复核。序列顶部的普通人形表示组织外部参与者，圈内人形表示组织岗位；改进图以斜线圆表示承担业务职责的软件系统。岗位自调用表示人员查阅或登记资料。分析重点在于三处信息衔接：取消与未到诊能否区分，临床清单移交后是否重复录入，以及季度收入与费用能否追溯来源。

\figoverview{02A_现状业务序列全景}{现状业务序列全景：单次就诊与季度经营复核}{fig:as-is-overview}

\figvectorlong{02_现状_预约与到诊}{现状业务序列：电话预约}{fig:as-is-booking}
前台先查找本店档案，再与牙医协调接诊时间，并向患者确认安排。人工确认是否同时涵盖牙椅、跨店任职及项目资质，是后续门店调研需要核实的内容；本模型暂不假定门店已具备统一的资源核验记录。
\figvectorlong{02B_现状_患者到诊}{现状业务序列：预约日到诊}{fig:as-is-arrival}
报到时，前台再次核对身份并询问资料变化，再将相关信息交给牙医。患者声明与牙医的临床风险判断分别由对应岗位记录，改进方案据此增加明确的到诊交接记录。
\figvectorlong{03_现状_诊疗与结算}{现状业务序列：临床诊疗}{fig:as-is-clinical}
诊疗环节形成检查所见与实际治疗记录。实施项目与最初预约诉求不一致时，应保留调整原因，并以牙医确认的实施事实作为后续结算依据。
\figvectorlong{03B_现状_诊疗费用结算}{现状业务序列：诊疗费用结算}{fig:as-is-billing}
收费员根据牙医移交的治疗清单核对数量和价格。手工重录存在漏项、重复及来源失真的风险，改进方案因而将结算信息核对与临床记录修改分属不同权限。
\figcompact{04_现状_季度核算}{现状业务序列：季度经营核算}{fig:as-is-quarter}
门店提交的收入、用品成本和运营费用若分类不同，总部即使汇总金额也无法可靠比较。尤其一位牙医跨店服务时，收入与费用需要按就诊和归属规则连接，而不能简单从某家门店总额中推算。

现状分析将问题归为身份、资源与核算三类：缺少共同患者标识会增加重复建档风险；任职、证书和排班分散时，预约条件难以一次核实；牙医跨店服务又使门店汇总不足以支持个人贡献核算。上述风险的发生频率及影响程度，仍须通过门店观察建立基线。

\subsection{改进后的岗位协作}
改进流程将资料检索、资源校验、记录保存与核算交由系统处理，身份确认、临床判断和财务审核仍由相应岗位负责。候选时段仅用于协商，预约提交时须重新核验资源；到诊登记记录报到事实，临床风险另由牙医复核。

\figoverview{05A_改进业务序列全景}{改进后业务序列全景：单次就诊与季度经营复核}{fig:to-be-overview}

\figvectorlong{05_改进后_预约与到诊}{改进后业务序列：电话预约}{fig:to-be-booking}
系统按诊所营业、牙医任职与资质、排班及牙椅占用情况提供候选时段，提交时再次核验并保存预约编号。前台据此向患者确认诊所、牙医与服务时间，形成资源条件和患者意愿均已核实的安排。
\figvectorlong{05B_改进后_患者到诊}{改进后业务序列：预约日到诊}{fig:to-be-arrival}
前台登记到诊并记录病史变化声明，系统将待复核信息交至牙医工作区。牙医完成本次风险复核后继续诊疗；两类确认分别保留责任人和时间，以便追溯交接过程。
\figvectorlong{06_改进后_诊疗与结算}{改进后业务序列：临床诊疗}{fig:to-be-clinical}
牙医依次复核风险、记录牙位观察和实际治疗，签署后提出完成请求。系统逐项核对完成条件，缺失时保留进行中状态并指出具体项目，避免结算环节才发现病历不完整。
\figvectorlong{06B_改进后_诊疗费用结算}{改进后业务序列：诊疗费用结算}{fig:to-be-billing}
结算行来自已完成就诊的实际治疗。收费员确认价格与数量后形成单据，收款与退款各保存独立流水。重复开单返回原编号，退款不抹去原收款，使患者和财务都能追查余额变化。
\figvectorlong{07_改进后_季度核算}{改进后业务序列：季度核对与关闭}{fig:to-be-quarter}
季度关闭包括输入核对、异常处理、结果复算与版本确认。费用仍有缺项或核算期间输入发生变化时，期间保持开放；核定快照锁定后供经营分析读取，后续更正通过授权重开处理。

由此，实际治疗、逐项应收、资金流水与期间结果构成连续且可追溯的记录链。收费员对来源有疑问时退回临床核查；财务在异常解决后重新计算并核定季度结果。经营查询分别呈现实时账本与已关闭期间快照，保持资金余额和核定利润的口径区别。

\subsection{交接问题与改进依据}
表~\ref{tab:handoffs} 将现状假设中的交接风险对应到共同记录与待确认制度。每处交接以可检查的业务结果界定；涉及机构政策的事项在门店确认后再形成强制规则。

\begin{longtable}{@{}L{23mm}L{45mm}L{44mm}L{31mm}@{}}
\caption{主要交接点的分析与处理}\label{tab:handoffs}\\\toprule
交接 & 现状风险 & 改进后的共同记录 & 上线前确认\\\midrule\endfirsthead
\toprule 交接 & 现状风险 & 改进后的共同记录 & 上线前确认\\\midrule\endhead
患者—前台 & 不同门店可能按电话或姓名重复建档 & 全网患者编号、重复候选及人工复核结果 & 身份核验材料\\
前台—牙医 & 预约时未见病史变化与治疗资格 & 到诊声明、牙医风险复核、服务日资质 & 复核责任分工\\
牙医—收费员 & 口头或纸质治疗清单重复录入 & 已实施治疗行作为结算来源 & 价格审核权限\\
收费员—财务 & 实时收款与季度已确认收入混淆 & 独立的收款流水与期间核定快照 & 跨期退款口径\\
诊所—总部 & 门店费用分类不一，比较失真 & 统一费用类别、归属期间及凭据 & 共用费用分摊\\
总部—门店 & 资质档案更新后旧预约可能失效 & 带生效日期的任职和证书记录 & 证书核验周期\\
\bottomrule\end{longtable}

改进流程的预期收益有明确测量对象：跨店复诊时需重新登记的字段数，预约因资源或资质冲突被改动的次数，临床明细与收费明细的差异笔数，季度关闭前未归属费用的数量。只有试点前后使用相同口径取数，才能评价系统是否实际降低工作量。

\subsection{预约结果的业务口径}
预约取消与未到诊可能使原定诊疗未能实施，但其发起原因、通知时点及资源能否再次利用并不相同。系统分别记录发起方、发生时间、资源释放与到诊情况，报表再依据明确的事件定义汇总，避免用单一状态掩盖不同的经营影响。

\begin{table}[!htbp]\centering\small
\caption{预约相关事件的区分}\label{tab:appointment-events}
\begin{tabular}{@{}L{27mm}L{41mm}L{41mm}L{36mm}@{}}\toprule
事件 & 成立条件 & 保存的关键事实 & 指标处理\\\midrule
改期 & 服务前更换有效时段 & 原时段、新时段、原因与执行人 & 不误记为一次新患者预约\\
主动取消 & 服务前明确取消 & 发起方、时间、原因、是否可补位 & 与未到诊分开统计\\
未到诊 & 到诊截止时间仍无报到 & 原预约时段、确认方式及记录人 & 单独计入爽约口径\\
到诊后未治疗 & 已报到，治疗未发生 & 到诊事实、牙医判断或其他原因 & 不归入未到诊\\
候补补位 & 释放时段被候补患者确认 & 原释放事件、邀约、确认及新预约 & 统计资源挽回，不抹去原取消\\
\bottomrule\end{tabular}\end{table}

例如“取消率”的分母可选已确认预约数，分子可只取服务日前的主动取消；“未到诊率”另取应到诊预约作分母。具体截止时间由机构制定，系统不预设惩罚或费用。候补成功后，原取消和新预约都保留，避免用补位覆盖原始事件而使取消率看似降低。

\subsection{从案例问题到可核实数据}
案例提出跨店诊疗标准、取消预约、牙医专业配置和季度贡献四项管理关切。前两项需要日常业务数据，后两项还依赖主数据质量和财务制度。表~\ref{tab:case-evidence} 将问题拆成观察量及其局限，避免把单一指标当成管理结论。

\begin{table}[!htbp]\centering\small
\caption{案例问题与数据依据}\label{tab:case-evidence}
\begin{tabular}{@{}L{30mm}L{57mm}L{58mm}@{}}\toprule
管理问题 & 可观察的数据 & 解释时的限制\\\midrule
跨店诊疗是否一致 & 同类治疗的实施项目、复诊与补记情况 & 患者病情不同，不能只比较治疗次数\\
取消带来多少损失 & 取消时点、牙椅空置、候补补位、未到诊 & 空闲时段不等于已发生收入损失\\
专业能力能否覆盖需求 & 有效资质、跨店任职、治疗需求与时段 & 证书有效不代表该时段牙医可排班\\
牙医季度贡献如何核算 & 实施收入、材料成本、费用归属与退款 & 未分摊费用时只能给出不完整估计\\
\bottomrule\end{tabular}\end{table}

经营页面因此须同时显示比率的分子、分母、日期与组织范围，并说明成本缺项、样本不足和指标不适用等情况。课程系统用于验证数据与计算关系，经营结论则需结合门店背景，由负责人和财务共同审阅。

\FloatBarrier
\figoverview{08_系统用例全景}{安德鲁牙科连锁运营管理系统用例全景图}{fig:system-usecase-overview}
\section{B 系统需求}
\subsection{系统边界与用例}
ADC Nexus 的执行者包括直接使用软件的人员及提供协助的外部身份服务。患者经本人入口读取已绑定资料，短信和微信平台仅参与身份验证。图形沿用课程参考的业务符号：人形头部和用例椭圆带斜线，外部软件采用非人业务实体符号并注明名称。指向用例的实线箭头表示角色参与该目标；处理次序、访问条件及例外路径由规约说明\cite{omg}。

系统用例全景按患者与预约、临床诊疗、诊疗结算、人员资源、经营治理和授权入口编排 18 项用户目标。为保证总览可读，微信与短信的辅助关联在身份验证分图中展开；同一“诊所经理”符号因版面需要重复绘出，表示同一岗位，并非两个角色。全景表示目标范围，不等于全部功能已验收。
\figuc{09_系统用例_患者与预约}{患者与预约系统用例图}{fig:uc-patient}
\figuc{10_系统用例_临床与结算}{临床与结算系统用例图}{fig:uc-care}
\figuc{11_系统用例_网络经营与治理}{经营与组织管理系统用例图}{fig:uc-management}
\figuc{11B_系统用例_财务与账户}{财务与账户管理系统用例图}{fig:uc-account}
\figuc{11C_系统用例_排班与牙椅}{排班与牙椅资源管理系统用例图}{fig:uc-resources}
\figuc{21_系统用例_门户与身份}{患者查询与身份认证系统用例图}{fig:uc-identity}

\begin{longtable}{@{}C{15mm}L{47mm}L{77mm}@{}}\caption{系统用例目录}\label{tab:uc-list}\\
\toprule 编号 & 用例 & 主执行者或业务对象\\\midrule\endfirsthead
\toprule 编号 & 用例 & 主执行者或业务对象\\\midrule\endhead
UC-01 & 建立并维护患者档案 & 前台；跨店患者身份\\
UC-02 & 安排诊疗预约 & 前台；牙医、诊所和牙椅资源\\
UC-03 & 变更诊疗预约 & 前台；改期、取消及历史\\
UC-04 & 办理患者到诊 & 前台；到诊事实与临床交接\\
UC-05 & 安排候补预约 & 前台；释放时段与患者确认\\
UC-06 & 记录临床就诊 & 牙医；风险、观察与实际治疗\\
UC-07 & 制定治疗计划 & 牙医；方案与患者决定\\
UC-08 & 完成临床就诊 & 牙医；记录完整性与签署\\
UC-09 & 生成诊疗结算单 & 收费员；实际治疗的逐项应收\\
UC-10 & 处理收款与退款 & 收费员；资金事实和余额\\
UC-11 & 维护牙医资质与任职 & 总部运营；服务日能力\\
UC-12 & 关闭季度财务期间 & 总部财务；收入、费用和结果版本\\
UC-13 & 分析经营绩效与资质覆盖 & 管理岗位；门店与全网比较\\
UC-14 & 管理账户与访问范围 & 系统管理员；岗位与诊所授权\\
UC-15 & 查阅本人诊疗信息 & 患者；预约、就诊摘要和账单\\
UC-16 & 验证登录身份 & 用户；密码、短信或微信身份\\
UC-17 & 管理账号密码 & 用户；设置、更改和找回密码\\
UC-18 & 维护排班与牙椅资源 & 诊所经理；牙医班次、停诊例外和牙椅可用性\\
\bottomrule\end{longtable}

业务用例与系统用例采用不同的观察层次。BUC-01 由建档、到诊和临床记录等多个系统任务共同支持；季度核算则属于机构内部管理需求。下文对全部 18 项系统用例分别给出规约。

\subsection{岗位权限与质量要求}
表~\ref{tab:permissions} 列出岗位的默认职责。实际访问还须结合管理员授予的诊所范围和生效期：总部岗位可读取授权统计，跨店明细则另行核验敏感资料权限。

\begin{table}[!htbp]\centering\small
\caption{岗位与主要业务权限}\label{tab:permissions}
\begin{tabular}{@{}L{24mm}C{17mm}C{17mm}C{17mm}C{17mm}C{17mm}C{17mm}@{}}\toprule
岗位 & 患者建档 & 预约到诊 & 临床签署 & 结算收退 & 本店分析 & 期间关闭\\\midrule
前台 & 写 & 写 & — & — & 读 & —\\
牙医 & 读 & 读 & 写 & 读 & 读 & —\\
收费员 & 读 & 读 & 读 & 写 & 读 & —\\
诊所经理 & 读 & 读 & 读 & 读 & 读 & —\\
总部财务 & 脱敏 & — & — & 汇总 & 汇总 & 写\\
总部运营 & 受限 & 汇总 & — & — & 汇总 & —\\
患者 & 本人 & 本人 & 摘要 & 本人 & — & —\\
\bottomrule\end{tabular}\end{table}

“读”仍受诊所范围、治疗关系及必要性限制；“本人”必须经账户绑定。管理岗位并不因职级较高而获得临床记录编辑权。质量验收首先看正确性和可追溯性：冲突预约不能成立，结算来源可查，拒绝请求不改写原记录。其次看可用性：前台应能在一个任务路径中完成查找、预约和到诊，牙医能在临床页面同时看见过敏提醒与历史摘要。性能与可靠性指标需要在符合目标诊所规模的数据集上测试，不能由开发机上的短时响应推定。

\section{B 系统用例规约}
规约采用同一粒度：执行者、前置条件、成功结果、正常处理以及会改变结果的异常。身份与权限由服务端检查；在页面上隐藏入口不能代替访问控制。性能目标和当前实现范围见第~\ref{sec:implementation}~节。

\spec{UC-01}{建立并维护患者档案}
{前台已登录并具有相应诊所的建档权限；取得患者提供的最小身份线索。}
{保存全网唯一的患者编号，并按本次确认结果记录联系方式、紧急联系人、过敏及病史变化。}
{1. 按姓名、出生日期和电话查找候选；2. 前台核对身份，选取既有档案或新建；3. 录入资料并提交；4. 系统检查重复风险和必填项，保存结果及变更记录。}
{疑似同人时先人工复核，不自动合并；“未提供病史”和“无已知过敏”不能混用。修改联系方式不覆盖已签署的临床记录。}

\spec{UC-02}{安排诊疗预约}
{前台有预约权限；患者身份、目标治疗、诊所和期望时间已确定。}
{形成带编号的预约，保存患者、牙医、时段及初始状态历史；启用牙椅排程时记录指定牙椅。}
{1. 输入治疗与时间条件；2. 系统返回可行牙医和时段；3. 前台与患者确认；4. 提交方案；5. 系统再次校验并占用资源，返回预约编号。}
{提交时发现患者或牙医冲突，或任职、资质、营业时间不符时，拒绝保存并给出原因；指定牙椅时还须检查牙椅冲突。候选查询不是占用承诺。}

\spec{UC-03}{变更诊疗预约}
{前台有该诊所预约操作权；预约尚未进入不允许变更的阶段。}
{保留原预约编号和历史，更新可执行时段；取消时记录发起方、时间与原因。}
{1. 读取预约及版本；2. 选择改期或取消；3. 提交新条件或理由；4. 系统校验状态与资源，保存新版本并释放旧占用。}
{已形成临床就诊的预约不得借取消清除诊疗事实。取消与未到诊使用不同状态；版本变化或新时段冲突时不覆盖他人操作。}

\spec{UC-04}{办理患者到诊}
{前台有当日预约的诊所范围；患者到达并可核对身份。}
{预约记录到诊时间，必要的资料变化声明进入牙医待复核事项。}
{1. 查找当日预约；2. 核对患者身份和联系方式、病史变化；3. 提交到诊；4. 系统登记并向临床工作区提供待诊资料。}
{身份不符时不得返回其他患者的临床资料。重复到诊不建立第二个有效就诊；前台登记声明不等于牙医确认风险。}

\spec{UC-05}{安排候补预约}
{前台有预约权限；患者接受候补，且已有释放或可供选择的时段。}
{患者确认后形成正式预约，候补条目和邀请结果均可追查。}
{1. 录入地点、时间及治疗限制；2. 系统按必需条件筛选；3. 前台联系候选患者；4. 患者确认后系统复核资源并建立预约。}
{拒绝、超时或资源变化均不产生正式预约。同一时段不能同时向多人保留有效确认窗口。当前程序的邀约确认写入尚未完成。}

\spec{UC-06}{记录临床就诊}
{负责牙医已登录；患者已到诊并建立与预约关联的临床就诊。}
{本次症状、检查、风险处理、观察、实际治疗和记录版本可追溯。}
{1. 牙医查看既往资料并复核过敏；2. 记录检查及牙位观察；3. 登记实际治疗、数量、牙位和依据；4. 保存记录并签署。}
{未核实的严重风险保留待处理状态。实际治疗可与预约诉求不同，但追加项目须说明原因；修改已签署记录必须追加版本。}

\spec{UC-07}{制定治疗计划}
{牙医已取得本次检查资料并具有相应临床权限。}
{保存方案版本、计划项目及患者对各项目的接受、拒绝或暂缓决定。}
{1. 牙医录入方案、费用说明和风险；2. 向患者解释；3. 记录患者逐项决定；4. 系统保存计划及必要的同意依据。}
{计划项目不直接计费，方案变更形成新版本。必要同意证据缺失时列为待核对并阻断就诊完成；已发生的治疗仍须如实记录。完整计划写入仍属待完善范围。}

\spec{UC-08}{完成临床就诊}
{牙医负责该次就诊；临床记录已签署。}
{就诊和关联预约记录完成事实及时间，治疗资料可供后续结算。}
{1. 牙医请求完成；2. 系统检查风险、观察、实际治疗、签署与服务日资格；3. 通过后在同一事务中保存就诊和预约的完成状态。}
{缺项时逐项返回，原状态不变。完成后的补记保留原事实、操作者和理由，不重新改写完成时间。}

\spec{UC-09}{生成诊疗结算单}
{收费员有该诊所权限；临床就诊已完成且可结算。}
{一张结算单及其逐项明细关联该就诊，保留项目名称、单位、数量、单价和金额快照。}
{1. 收费员读取完成就诊；2. 系统取得实际治疗及价格依据；3. 核对各行并计算总额；4. 保存结算单并返回编号。}
{临床来源不完整时停止开单；同一就诊重复请求返回既有单据，不生成第二张。诊疗结算单不是法定税务发票。}

\spec{UC-10}{处理收款与退款}
{收费员有财务权限；结算单有效，退款有原收款与授权理由。}
{收款、退款分别保存凭据，单据余额与状态按有效净收款更新。}
{1. 展示应收和已收；2. 收费员核实付款凭据并登记金额；3. 退款时选择原收款、填写理由；4. 系统校验后保存交易并重算余额。}
{收款不得超过余额，累计退款不得超过可退原收款；失败不修改财务记录。系统登记已核实的资金事实，不执行银行或微信支付划转。}

\spec{UC-11}{维护牙医资质与任职}
{总部运营人员取得人员主数据维护权限。}
{保存牙医的多项证书及跨店任职、生效期和变更依据。}
{1. 定位牙医；2. 录入资质证明、颁发方及有效期；3. 设置诊所任职和工作时间；4. 系统校验并保存历史。}
{预约按服务日、治疗按实施日判断有效性。失效证书保留历史但不用于未来资格匹配；证件核验与审批工作流仍待完善。}

\spec{UC-12}{关闭季度财务期间}
{总部财务有期间关闭权限；收入、用品成本和运营费用的归属资料已进入核对范围。}
{保存核定结果和输入版本，锁定期间，并根据批准口径确定全网激励基金总额。}
{1. 财务发起检查；2. 系统列出缺项和差异；3. 门店处理、财务复核后重算；4. 财务核对结果版本并关闭期间。}
{存在阻断事项或版本已变化时不得关闭。普通操作不能改写锁定结果；重开须授权和理由。牙医费用分摊与完整贡献计算尚未验收。}

\spec{UC-13}{分析经营绩效与资质覆盖}
{总部或诊所管理人员已登录，并具有相应组织范围。}
{在指定日期与诊所范围内取得预约履约、治疗结构、财务和有效资质比较结果。}
{1. 选择日期、诊所和指标；2. 系统按授权范围计算并标注分子、分母及口径；3. 用户查看异常来源或扩张情景。}
{分母为零时显示“不适用”。诊所经理不得从统计总量推知其他店敏感数据；情景分析不直接改变实际任职或预约。}

\spec{UC-14}{管理账户与访问范围}
{系统管理员已具备账户管理权限。}
{账户、角色及诊所范围按生效期保存，停用或重置后相关会话失效。}
{1. 定位人员账户；2. 设置岗位、诊所和有效期；3. 系统校验授权冲突；4. 保存分配与审计事件。}
{管理员身份不默认授予临床或财务操作权。接口对列表、详情和修改执行相同范围检查，拒绝请求不泄露无权记录。}

\spec{UC-15}{查阅本人诊疗信息}
{患者完成身份验证，账户已由机构核实并绑定患者编号。}
{只返回本人预约、就诊摘要和诊疗结算单。}
{1. 患者进入本人入口；2. 系统按账户绑定读取信息；3. 患者查看详情及必要说明。}
{没有绑定关系时引导机构核实，不能通过输入任意患者编号建立绑定。临床内部备注和他人资料不向患者入口开放。}

\spec{UC-16}{验证登录身份}
{用户持有有效账户或符合机构开户条件。}
{核验成功后创建受权限与诊所范围约束的会话。}
{1. 用户选择账号密码、手机验证码或微信身份；2. 系统完成相应核验；3. 读取账户状态及授权；4. 建立会话并返回允许访问的入口。}
{验证码短期有效且一次消费；密码尝试限频。微信回调须核验状态和绑定。正式短信、微信服务缺少机构凭据时明确失败，不用模拟验证码替代。}

\spec{UC-17}{管理账号密码}
{已登录用户可设置或更改密码；未登录找回者须通过已绑定手机号核验。}
{只更新本人密码摘要，重置后旧会话失效。}
{1. 用户选择设置、更改或找回；2. 系统核验原密码或手机挑战；3. 用户输入新密码；4. 系统检查策略并保存摘要。}
{手机验证失败或手机号未绑定时不改变凭据。界面不回显原密码，日志不写入明文密码或验证码。}

\spec{UC-18}{维护排班与牙椅资源}
{诊所经理已获本店资源维护权；待调整牙医的任职和目标诊所均已确认。}
{牙医班次、停诊例外或牙椅可用状态形成生效版本；受影响预约被列明，不被静默改写。}
{1. 经理选择服务日期、牙医或牙椅及拟变更时段；2. 系统读取任职、排班、牙椅状态与现有预约；3. 显示冲突和受影响患者；4. 经理处理例外并确认；5. 系统按版本保存资源变化与操作依据。}
{与已确认预约冲突、任职无效或版本过期时不得直接生效，须先协调预约或重新提交。单独修改牙医班次或牙椅状态时，仅更新对应资源；资源管理审批制度仍需门店确认，列为目标设计而非现有程序已验收能力。}

\subsection{关键用例的步骤与校验}
上述规约用于确定业务目标；下列核对表用于评审跨对象的提交点。表中“提交”是一次持久化边界，而非按钮名称。对需要人工确认的事项，系统只能保存确认结果及依据，不能自行作出临床或财务判断。

\begin{longtable}{@{}C{12mm}L{46mm}L{51mm}L{41mm}@{}}
\caption{UC-02 安排预约：逐步校验}\label{tab:booking-check}\\\toprule
步骤 & 输入或动作 & 系统检查 & 留存结果\\\midrule\endfirsthead
\toprule 步骤 & 输入或动作 & 系统检查 & 留存结果\\\midrule\endhead
1 & 前台定位患者 & 诊所范围、患者身份和重复线索 & 已确认的患者编号\\
2 & 选择治疗需求、诊所、日期 & 目录是否适用、诊所是否营业 & 候选查询条件\\
3 & 查看牙医与牙椅候选 & 任职、证书在服务日是否有效；时段是否重叠 & 仅显示候选，不锁定资源\\
4 & 与患者确认具体时段 & 必要联系方法是否可用 & 确认内容与操作者\\
5 & 提交预约 & 在写入时重做资格和占用检查 & 预约编号、初始状态历史\\
6 & 资源已变化 & 不生成预约；返回变化原因 & 保留输入，提示重选\\
\bottomrule\end{longtable}

预约的“可选”与“已占用”有时间差。若同时提交的请求只各自读取到空闲，再分别写入，页面层的可选提示不能防止双占。因而上线验收除普通冲突场景外，还必须对同一牙医、同一牙椅的并发请求分别验证。

\begin{longtable}{@{}C{12mm}L{45mm}L{52mm}L{41mm}@{}}
\caption{UC-06／08 临床记录与完成：逐步校验}\label{tab:clinical-check}\\\toprule
步骤 & 牙医操作 & 提交前检查 & 保存事实\\\midrule\endfirsthead
\toprule 步骤 & 牙医操作 & 提交前检查 & 保存事实\\\midrule\endhead
1 & 打开已到诊的患者 & 负责牙医、诊所、就诊状态 & 访问及读取范围\\
2 & 查阅并复核过敏、病史 & 未知、无已知与已知风险分开呈现 & 本次复核人与时间\\
3 & 记录主诉和牙位检查 & 牙位编码及必要描述 & 检查观察\\
4 & 记录实际治疗 & 项目资格、数量、牙位及计划偏离理由 & 实施项目行\\
5 & 保存或修订临床记录 & 版本号及签署状态 & 新版本，不覆盖旧版本\\
6 & 请求完成就诊 & 风险处理、必要同意、记录签署、实施项完整性 & 完成时间或逐项阻断原因\\
\bottomrule\end{longtable}

计划是向患者解释后的方案，实际治疗是本次已经执行的事实。患者暂缓某项方案时，计划仍可保留，但不得由此产生收费行。完成后的补记应以带理由的新版本关联原就诊，不能让历史签署内容无痕变化。

\begin{longtable}{@{}C{12mm}L{46mm}L{51mm}L{41mm}@{}}
\caption{UC-09／10 结算与收退款：逐步校验}\label{tab:billing-check}\\\toprule
步骤 & 收费员操作 & 系统检查 & 保存结果\\\midrule\endfirsthead
\toprule 步骤 & 收费员操作 & 系统检查 & 保存结果\\\midrule\endhead
1 & 读取已完成就诊 & 是否属于授权诊所、是否已有有效单据 & 就诊来源编号\\
2 & 核对实际治疗与价格 & 项目、数量、价格快照、舍入规则 & 结算行与应收总额\\
3 & 确认生成 & 同一就诊唯一性、版本未变化 & 单据编号与创建时间\\
4 & 登记已核实收款 & 金额不超过余额，凭据未重复 & 收款流水及未收余额\\
5 & 选择原流水办理退款 & 退款权限、原款余额和理由 & 独立退款流水\\
6 & 查看单据状态 & 净收款与应收金额一致 & 部分收款、结清或退款后状态\\
\bottomrule\end{longtable}

诊疗结算单与税务发票不同；现金、银行卡或第三方支付是否实际成功，应由机构收款凭据证实。退款也不是直接把原收款金额改小，否则无法解释季度内的收退款时间差。

\begin{longtable}{@{}C{12mm}L{46mm}L{51mm}L{41mm}@{}}
\caption{UC-12 季度期间关闭：逐步校验}\label{tab:period-check}\\\toprule
步骤 & 财务操作 & 系统检查 & 保存结果\\\midrule\endfirsthead
\toprule 步骤 & 财务操作 & 系统检查 & 保存结果\\\midrule\endhead
1 & 选择待关闭季度 & 期间是否开放、操作人是否有总部权限 & 待核算范围\\
2 & 检查诊所收入 & 未结项目、退款及跨期事项 & 收入清单与差异\\
3 & 核对成本及运营费用 & 凭据、归属诊所、期间与分摊方法 & 缺项或待复核清单\\
4 & 重算并解释异常 & 输入版本是否一致、核算口径是否批准 & 核算快照与计算依据\\
5 & 确认关闭 & 阻断项是否清零、输入是否被并发修改 & 锁定时间、核定人和版本\\
6 & 申请重开 & 授权、理由与后续影响 & 重开事件和复核任务\\
\bottomrule\end{longtable}

期间关闭不应自动清除异常。对于不能分配到牙医的共用费用，可先形成门店层的待归属项，待财务确认分摊规则后再计算个人贡献；若规则未批准，基金与贡献结果应标明口径不完整。

\begin{longtable}{@{}C{12mm}L{46mm}L{51mm}L{41mm}@{}}
\caption{UC-16／17 身份验证与密码找回：逐步校验}\label{tab:auth-check}\\\toprule
步骤 & 用户动作 & 系统检查 & 失败时处理\\\midrule\endfirsthead
\toprule 步骤 & 用户动作 & 系统检查 & 失败时处理\\\midrule\endhead
1 & 选择登录方式 & 手机、密码或微信入口的可用配置 & 未配置方式明确提示\\
2 & 请求短信验证码 & 号码格式、限频及供应商返回 & 不提供本地模拟码\\
3 & 提交验证码或密码 & 时效、一次消费、密码摘要比较 & 统一错误提示并限频\\
4 & 完成微信回调 & 回调状态、签名及已绑定身份 & 不据未核实身份建会话\\
5 & 建立会话 & 账户状态、角色及诊所范围 & 账号停用时拒绝\\
6 & 找回或修改密码 & 已绑定手机号挑战或原密码、新密码策略 & 不改变原凭据\\
7 & 保存新密码 & 单向密码摘要、版本及旧会话撤销 & 旧会话不可继续访问\\
\bottomrule\end{longtable}

\subsection{需求追踪与评审结果}
需求追踪同时覆盖成功操作与拒绝路径。G1 对应 UC-01、04、15，核对重复档案和本人绑定；G2 对应 UC-02、03、11、18，核对服务日资质、排班与牙椅冲突；G3 对应 UC-06 至 UC-10，核对签署、实际治疗、结算来源和余额；G4 对应 UC-12、13，核对费用归属、期间快照与范围授权。UC-14、16、17 提供共同的身份和权限前提。尚未满足的规则继续作为设计要求，在验收表中分别标注待专项验收或待开发验收。

\Needspace{7\baselineskip}
\section{C 系统分析}
\subsection{设计原则与职责划分}
患者身份、临床事实与资金事实分别由相应岗位维护：前台核实身份并管理预约，牙医记录检查与治疗，收费员据实际治疗生成结算单，总部财务核定期间结果。系统围绕职责边界检查记录一致性并保存修改依据；同一对象的更新带有版本条件，防止旧页面提交覆盖新记录。

服务端按“账户状态—岗位权限—诊所范围—记录归属”顺序判断访问。患者本人查询另以已核实的账户—患者绑定为条件；知道患者编号不构成访问依据。跨店复诊允许有权限的医护人员读取必要历史，但不能因此把所有门店的敏感资料开放给普通门店账号。

\subsection{领域对象与关系}
领域类图按患者身份、预约资源、临床事实与财务事实分组，类框依次列示名称、属性和操作。普通关联用实线，单向导航用开放箭头；聚合与组合的空心、实心菱形均置于整体端，重数注明于对应关联端。泛化以指向父类的空心三角表示。中文名称表达业务含义，String、DateTime、Decimal 等为逻辑类型，具体存储类型在数据层确定。图中关系与表~\ref{tab:domain-cardinality}、相邻正文共同说明数量和生命周期约束。

患者以全网稳定编号识别，电话号码作为可变联系方式保存。过敏、同意材料及既往就诊均记录来源和时间；疑似重复档案经授权人员复核后处理，合并后仍须能够追溯原编号。

\figclassvector{12_领域类图_患者与预约}{患者与预约领域类图}{fig:domain-patient}
\figclassvector{12E_领域类图_患者身份}{患者身份、过敏与绑定领域类图}{fig:domain-identity}

\Needspace{7\baselineskip}
预约必须指向患者、诊所、牙医及服务时间；牙椅关联允许为空，分配牙椅后才进行牙椅占用校验。同一牙医不能在重叠时段接受两个有效预约，已分配的同一牙椅也不能重复占用。资质与任职按服务日判断，而非只查看当前状态。取消、未到诊与改期分别记录原因、操作者和发生时间，供后续分析资源损失。预约成立并不自动表示患者已到诊。

\figclassvector{13_领域类图_人员资质与目录}{牙医任职、资质与治疗目录领域类图}{fig:domain-workforce}

一次就诊关联预约但有独立状态。患者诉求、检查观察、拟议治疗和实际治疗是不同事实：计划可以调整，账单只引用已实施且经核对的项目。临床记录保留版本、作者与签署时间；过敏复核、必要的同意证据和签署要求属于完成就诊的阻断条件。该设计遵循患者病史及时更新和诊疗记录可追溯的基本要求\cite{ada-history,ada-records}。

\figclassvector{12B_领域类图_临床诊疗}{就诊与实际治疗领域类图}{fig:domain-clinical}
\figclassvector{12D_领域类图_记录与风险}{临床记录版本与过敏复核领域类图}{fig:domain-risk}
\figclassvector{12C_领域类图_计划与同意}{治疗计划、项目与同意证据领域类图}{fig:domain-plan}

计划项目表达拟议治疗，实际治疗表达已经实施的事实，两者不能互相代替。程序已有计划、实际治疗、临床记录版本和过敏复核等数据结构；图中的逐项目同意证据是拟完善的设计约束，不能仅凭患者的通用隐私同意推断某项治疗已获知情同意。完成门槛仍须在岗位试用和专项测试中核实。

\figclassvector{13B_领域类图_收费与季度财务}{结算单、明细与收退款领域类图}{fig:domain-finance}

诊疗结算按一次就诊建立唯一单据。结算行保存治疗项目、数量、单价及形成时的价格快照，目录日后调价不改写历史单据。收款和退款各自保留独立流水，净收款由流水计算；作废也保留原始单据及理由。季度核算区分门店经营结果与牙医个人贡献。案例提出按季度净利润提取 10\% 用于奖金\cite{case}，因此基金只在财务确认的全网净利润口径上计算，不在图中擅自规定个人奖金的分配公式。

\figclassvector{13C_领域类图_季度贡献}{季度结果与牙医贡献领域类图（目标设计）}{fig:domain-contribution}
\figclassvector{13D_领域类图_费用归属}{期间费用与归属领域类图（目标设计）}{fig:domain-allocation}

现有结构已独立保存期间结果，牙医级贡献与运营费用分摊仍属目标设计。分摊记录须关联原始费用、金额和经财务批准的依据，再据此计算个人贡献。图~\ref{fig:domain-contribution} 与图~\ref{fig:domain-allocation} 说明这一目标数据边界。

\begin{table}[!htbp]\centering\small
\caption{主要领域关联的数量约束}\label{tab:domain-cardinality}
\begin{tabular}{@{}L{38mm}L{106mm}@{}}\toprule
关联 & 约束\\ \midrule
患者、牙医、诊所与预约 & 每次预约各对应一名患者、一名牙医和一家诊所；三类主体均可关联多次预约\\
牙椅与预约 & 每次预约可尚未分配牙椅；分配后对应一个牙椅，且须满足时段不冲突\\
预约与临床就诊 & 一次预约至多形成一次就诊；未到诊、取消等情况可以不形成就诊\\
就诊与记录、实施项目 & 每条记录和实施项目归属于一次就诊；就诊可有多条记录及实施项目\\
就诊与诊疗结算单 & 一次就诊至多形成一张结算单；未完成诊疗时可以尚未开单\\
实际治疗与结算明细 & 每条明细有确定的实施项目来源；同一实施项目至多计费一次\\
结算单、收款与退款 & 一张单据可追加多笔收款；每笔退款指向一笔原收款，累计退款不超过可退金额\\
患者与账户绑定 & 患者及本人账户各至多建立一个有效绑定，须经身份核实\\
\bottomrule\end{tabular}\end{table}

\subsection{状态与业务约束}
四张状态图分别描述预约、临床就诊、诊疗结算单及退款后的结算状态。图中 T 标识变迁，邻接表列出触发事件与守卫条件。改期保留预约编号并重新校验资源；完成诊疗须满足风险复核与签署要求；结算状态根据净收款重算，原收款及退款流水均予保留。已结清仍可因退款转出，因此不标为不可逆的终态。

\figstatevector{14_状态机_预约}{预约状态机}{fig:state-appointment}
\begin{table}[H]\centering\small
\caption{预约状态变迁与守卫条件}\label{tab:state-appointment}
\begin{tabular}{@{}C{10mm}L{45mm}L{90mm}@{}}\toprule
编号 & 触发事件 & 主要条件\\\midrule
T1 & 创建预约 & 患者、牙医、诊所、时段及必要资质校验通过\\
T2 & 登记到诊 & 身份与预约匹配，预约尚有效\\
T3 & 开始诊疗 & 到诊交接和待处理风险已核对\\
T4 & 完成就诊 & 实施项目、病历签署及必要同意材料已核对\\
T5 & 取消预约 & 尚未到诊，保留取消原因与操作人\\
T6 & 撤销到诊 & 尚未实施治疗，授权人员记录更正原因\\
T7 & 标记未到诊 & 已超过约定时点及保留时限，未办理到诊\\
\bottomrule\end{tabular}\end{table}

\clearpage
\figstatevector{15_状态机_临床就诊}{临床就诊状态机}{fig:state-visit}
\begin{table}[H]\centering\small
\caption{临床就诊状态变迁与守卫条件}\label{tab:state-visit}
\begin{tabular}{@{}C{10mm}L{43mm}L{92mm}@{}}\toprule
编号 & 触发事件 & 主要条件\\\midrule
T1 & 建立就诊 & 到诊及预约资料有效\\
T2 & 进入待诊疗 & 无待处理的过敏、同意或资质风险\\
T3 & 进入安全复核 & 存在需要人工确认的风险事项\\
T4 & 完成安全复核 & 复核结论、处理人及时间已留痕\\
T5 & 开始诊疗 & 牙医任职与项目资质有效\\
T6 & 完成就诊 & 实施项目和病历版本均已核对\\
\bottomrule\end{tabular}\end{table}

\clearpage
\figstatevector{16_状态机_发票}{诊疗结算单收款状态机}{fig:state-invoice}
\begin{table}[H]\centering\small
\caption{诊疗结算单状态变迁与守卫条件}\label{tab:state-invoice}
\begin{tabular}{@{}C{10mm}L{43mm}L{92mm}@{}}\toprule
编号 & 触发事件 & 主要条件\\\midrule
T1 & 生成结算单 & 同一就诊未重复开单，明细来自实际治疗\\
T2 & 部分收款 & 净收款大于零且小于应收金额\\
T3 & 全额收款 & 一次收款后净收款等于应收金额\\
T4 & 结清尾款 & 追加收款后净收款等于应收金额\\
T5 & 授权作废 & 净收款为零，记录作废原因与审批人\\
\bottomrule\end{tabular}\end{table}

\clearpage
\figstatevector{16B_状态机_退款}{退款后结算状态判定}{fig:state-refund}
\begin{table}[H]\centering\small
\caption{退款后的结算状态重新判定}\label{tab:state-refund}
\begin{tabular}{@{}C{10mm}L{43mm}L{92mm}@{}}\toprule
编号 & 触发事件 & 净收款判定\\\midrule
T1 & 已结清后部分退款 & $0 < \text{净收款} < \text{应收}$\\
T2 & 已结清后全额退款 & 净收款为零，单据仍保留\\
T3 & 部分收款后退至零 & 净收款为零\\
T4 & 部分收款后继续退款 & 退款后净收款仍大于零且小于应收\\
\bottomrule\end{tabular}\end{table}

\clearpage
\subsection{分析类与关键交互}
分析类划分为边界、控制与实体三类职责。图~\ref{fig:bce} 以安排预约为例：边界对象接收操作并呈现结果，控制对象核对访问范围并协调资源，实体对象维护占用与预约事实；虚线箭头表示依赖。随后六张分析序列图将关键用例分配到相关对象，进一步说明调用顺序、规则归属和提交边界。分析职责与实现文件或界面页面不必逐一对应。

\figvector{24_分析类_边界与控制}{安排预约的分析类协作}{fig:bce}

\begin{table}[H]\centering\small
\caption{分析类职责与交互边界}\label{tab:analysis-class-duty}
\begin{tabular}{@{}L{22mm}L{32mm}L{44mm}L{40mm}@{}}\toprule
交互场景 & 边界类职责 & 控制类职责 & 实体类职责\\ \midrule
安排预约 & 接收患者、时段和方案，呈现校验结果 & 核对诊所范围、任职资质与资源冲突，并控制一次提交边界 & 资源日历维护占用；预约及预约历史保存结果\\
完成临床就诊 & 提交就诊号和记录版本，展示缺项或完成结果 & 核对过敏复核、实际治疗、同意与签署完整性 & 临床就诊保存状态与版本；预约记录完成时间\\
生成诊疗结算单 & 按就诊号发起开单，展示逐项费用 & 防止重复开单，按实际治疗和价格快照建立明细 & 临床就诊提供计费来源；结算单保存逐项应收\\
\bottomrule\end{tabular}\end{table}

六张分析序列图依次覆盖预约安排、排班与牙椅维护、临床完成、结算开单、季度关闭和绩效查询。消息由发起对象指向承担职责的对象，虚线表示返回结果。预约在资源复核后保存主记录与历史，临床完成依据实际治疗和签署版本，结算生成则校验来源及重复开单；这些关键提交点在图中分别展开。宽幅图采用原尺寸矢量图版，以便核对对象名称与消息细节。

\figanalysisplate{18_分析序列_安排预约}{安排预约分析序列图}{fig:seq-booking}
\figanalysisplate{18B_分析序列_维护排班与牙椅}{维护排班与牙椅资源分析序列图}{fig:seq-resource-maint}
\figanalysisplate{20_分析序列_完成临床就诊}{完成临床就诊分析序列图}{fig:seq-clinical}
\figanalysisplate{19_分析序列_完成就诊并开票}{生成诊疗结算单分析序列图}{fig:seq-invoice}
\figanalysisplate{20C_分析序列_季度期间关闭}{季度财务期间关闭分析序列图（目标设计）}{fig:seq-quarter-close}
\figanalysisplate{20B_分析序列_经营绩效查询}{经营绩效与资源利用分析序列图}{fig:seq-performance}

排班与牙椅维护以联合调整为代表路径，单项变更仅写入对应实体。季度关闭先核对缺项、输入版本与候选结果，再在同一提交边界内锁定期间、保存核定快照及基金；版本变化时拒绝关闭。该协作属于目标设计，牙医级费用分摊口径仍待机构确认。绩效查询分别取得预约履约、牙医排班和牙椅占用的计算依据，核对实际治疗与结算来源，并读取季度快照。期间尚未关闭或费用待归属时，结果标为估算并列明缺项。

\section{D 系统设计}
\subsection{逻辑架构与边界}
系统通过浏览器提供员工工作台与患者本人入口。应用层组织身份、患者、预约、临床、结算及经营业务，数据层保存业务记录、版本和审计事件。短信与微信服务经适配接口接入；正式启用前须具备机构账号、回调域名和隐私告知等条件。配置缺失时明确提示通道不可用，不发放替代验证码。架构图中的外部平台表示目标集成边界。

\figvector{17_逻辑架构}{系统逻辑架构}{fig:architecture}

\subsection{部署边界与外部依赖}
图~\ref{fig:deployment} 区分终端、应用运行环境、持久化数据库与机构外部服务。浏览器通过加密连接访问应用；业务服务处理授权和事务，适配器负责已配置外部通道的请求与回执。课程实现以 Cloudflare 兼容环境及 SQLite/D1 为技术基线。正式部署还须核验域名、机构账号、供应商凭据和回调，并落实密钥保管与日志脱敏。

\figvector{25_部署边界图}{系统部署边界与外部依赖}{fig:deployment}

\subsection{界面导航与边界类映射}
预约工作台由日视图进入患者检索、时段选择与提交核对。发生冲突时保留已填资料，说明原因并引导调整时段。图~\ref{fig:ui-navigation} 描述界面导航，与前述预约对象的业务状态机分属不同层次；N1--N6 依次标识开始建约、选定患者、选择资源、提交成功、提示冲突和调整时段。表~\ref{tab:ui-boundary-map} 给出界面区域与分析边界职责的对应。

\figstate{23_预约界面导航状态}{预约工作台界面导航}{fig:ui-navigation}
\begin{table}[H]\centering\small
\caption{界面原型区域与分析边界类映射}\label{tab:ui-boundary-map}
\begin{tabular}{@{}L{29mm}L{33mm}L{82mm}@{}}\toprule
界面区域 & 对应边界职责 & 必要输入、反馈与限制\\\midrule
预约日视图 & 预约列表边界 & 按授权诊所、日期、牙医过滤；显示已占时段和状态\\
患者检索 & 患者选择边界 & 按稳定编号及必要线索检索；联系方式按岗位遮罩\\
时段与资源选择 & 资源选择边界 & 展示营业时间、停诊、牙医资质和牙椅可用性\\
提交核对 & 预约提交边界 & 汇总患者、项目、时段与告知信息；提交前再次校验\\
冲突与确认结果 & 反馈边界 & 区分版本冲突、资源冲突和成功；成功时返回预约编号\\
\bottomrule\end{tabular}\end{table}

\subsection{业务层：服务与提交边界}
业务服务按职责及一致性边界划分。患者服务维护全网身份、联系方式与风险资料，预约服务管理资源占用和状态历史，临床服务保存检查与实施事实，结算服务据已完成就诊形成应收并登记资金流水，期间服务汇总、核对及锁定季度结果。人员资质由预约、临床和经营分析共同引用；身份服务统一提供账户状态与访问范围判断。

\begin{longtable}{@{}L{26mm}L{49mm}L{52mm}L{22mm}@{}}
\caption{主要服务的输入与提交边界}\label{tab:services}\\\toprule
服务 & 主要输入 & 一次提交应保持一致的记录 & 调用者\\\midrule\endfirsthead
\toprule 服务 & 主要输入 & 一次提交应保持一致的记录 & 调用者\\\midrule\endhead
患者身份 & 身份线索、联系及过敏资料 & 患者主档与本次变更事件 & 前台、牙医\\
预约资源 & 患者、治疗、时段、牙医、牙椅 & 预约、状态历史与资源占用 & 前台\\
到诊交接 & 预约、核实结果、变更声明 & 到诊状态与待复核事项 & 前台\\
临床就诊 & 风险复核、观察、实施项目、记录 & 就诊版本、实施项和完成状态 & 牙医\\
诊疗结算 & 完成就诊、价格快照、收退款凭据 & 单据、行项目及独立资金流水 & 收费员\\
季度核算 & 门店收入、用品成本、运营费用 & 异常清单、输入版本与锁定快照 & 总部财务\\
权限治理 & 账户、岗位、诊所和有效期 & 授权分配与审计事件 & 管理员\\
\bottomrule\end{longtable}

一次提交中的关联写入应共同成功或撤回。预约主记录与状态历史、结算单与行项目均须保持完整。外部短信请求与数据库写入采用分步处理，避免等待供应商期间持有长事务：先保存挑战及发送状态，再依据供应商结果更新，发送失败时挑战不可用。微信回调须核验发起时保存的状态参数及身份绑定后，方可建立会话。

\subsection{持久化职责与数据访问顺序}
图~\ref{fig:persistence-class} 给出目标持久化职责：预约服务依赖仓储接口，D1 数据访问实现接口，并将预约、状态历史和牙椅资料映射到物理表。虚线表示依赖，空心三角表示接口实现；图中的模块分离属于设计方案。三个表对象分别对应 \texttt{appointments}、\texttt{appointment\_status\_history} 与 \texttt{clinic\_chairs}。

\figvector{21_持久化映射类图}{预约服务的持久化职责映射}{fig:persistence-class}

图~\ref{fig:data-access-sequence} 展示安排预约的目标数据访问路径。服务读取时段占用与版本，核对诊所范围、牙医任职、项目资质及资源冲突；拒绝分支返回具体原因，不写入预约。成功分支在同一事务中保存预约与状态历史，并于提交时复核版本和占用条件。键与唯一约束保障记录层面的完整性，时间区间重叠仍须由事务内的冲突检查及并发控制共同处理。

\figvectorlong{22_预约数据访问序列}{安排预约的数据访问序列}{fig:data-access-sequence}

\subsection{异常与恢复策略}
异常处理区分业务拒绝与系统故障。资质过期、牙椅冲突或余额不足时，返回相关字段及原因，供用户调整后重试；数据库不可用或供应商超时，则记录请求标识并提示后续处理。重复生成结算单按就诊编号返回既有结果，由服务端保证幂等性。预约、退款和期间关闭日志保留请求主体与结果，排除完整病历正文及验证码等敏感内容。

恢复测试应从备份重建临时环境，复查记录数量、金额平衡与最近状态，并核验预约—患者、临床—预约、结算—实际治疗及审计—修改事件四类关联。正式部署还须确定恢复点目标、备份保留期限、密钥管理与运维责任，形成可演练、可追责的恢复安排。

\subsection{关系数据库设计范围与建模口径}
数据层采用 SQLite/D1，物理结构以 Drizzle 迁移快照为依据。现有模型包含 61 张表、730 个字段、119 条已声明外键和 124 个索引。各表按生命周期与约束责任组织为九个业务簇：核心链覆盖患者识别、预约、就诊、收退款和季度核算，支撑域保存账户授权、通知、审计及扩张情景。联系人电话、证书号等可变属性与关系键分离，面向岗位的业务编号亦区别于稳定的内部主键。

\begin{table}[!htbp]\centering\small
\caption{关系数据库业务簇与规模}\label{tab:data-groups}
\begin{tabular}{@{}C{10mm}L{38mm}C{13mm}L{85mm}@{}}\toprule
序 & 业务簇 & 表数 & 主要责任\\\midrule
1 & 组织、账户与授权 & 10 & 组织、诊所、账户、角色权限、有效范围、系统配置与审计主体\\
2 & 身份认证与通知 & 6 & 外部身份映射、验证码挑战、会话、OAuth 状态、本人绑定及通知投递\\
3 & 患者主档与风险资料 & 6 & 患者、紧急联系人、过敏、同意、重复档案复核与时间线\\
4 & 牙医、资质与治疗目录 & 6 & 牙医任职、资质类型与有效期、治疗目录及项目资格要求\\
5 & 排班、资源与预约 & 10 & 营业时间、停诊、排班例外、牙椅、预约、候补与状态历史\\
6 & 临床就诊与病历版本 & 8 & 就诊、过敏复核、牙位观察、计划、实际治疗、记录版本与附件\\
7 & 结算、收款与成本 & 7 & 结算单、逐项明细、调整、独立收退款以及实际治疗成本\\
8 & 期间核算与激励 & 4 & 运营费用、关闭异常、期间快照和经批准的激励池\\
9 & 扩张情景与候选诊所 & 4 & 情景版本、候选诊所、可追溯假设和预测结果\\
\bottomrule\end{tabular}\end{table}

金额字段以人民币分保存，避免浮点舍入；时间点按 ISO 8601 存储并按诊所时区展示，业务日期与瞬时时间分开。状态历史、病历版本、收款、退款和审计事件采用追加式记录。结构化 JSON 只用于核算输入快照或尚不能稳定预定义的扩展载荷，可查询、可约束的重复事实仍使用关系表。

\subsection{全景与分域关系数据图}
关系数据图先呈现 61 张表的全景，再按九个业务簇展开。全景用于识别主链与跨域依赖，分域图用于核对字段和局部关系；十张图均采用可放大的原尺寸矢量 PDF，并附独立图版及可编辑源文件。浅色区分业务簇，实线箭头由外键表指向被引用表。全景保留跨域连线，分域图在外键字段行注明域外目标，使局部阅读仍能定位完整依赖。

表框按网格逐行列出物理字段。\texttt{PK} 表示主键，\texttt{FK} 表示已声明外键，\texttt{NN} 与 \texttt{NULL} 分别表示不可空和可空，\texttt{IDX} 表示参与现有索引，\texttt{DEF} 表示具有数据库默认值。全景、分域图、字段目录、关系清单和 DDL 由同一份迁移快照生成，采用一致的统计口径。大幅图版供电子阅读时放大核验；纸质阅读宜结合分域图及关系说明。

\begin{table}[H]\centering\small
\caption{关系数据图版阅读索引}\label{tab:erd-plates}
\begin{tabular}{@{}C{13mm}L{50mm}C{16mm}L{66mm}@{}}\toprule
图号 & 范围 & 表数 & 建议重点核对的连接\\\midrule
图~\ref{fig:erd-panorama} & 全景 & 61 & 患者—预约—就诊—结算与跨域授权\\
图~\ref{fig:erd-organization} & 组织、账户与授权 & 10 & 用户、角色、诊所范围、审计主体\\
图~\ref{fig:erd-authentication} & 身份认证与通知 & 6 & 会话、挑战、外部映射与投递记录\\
图~\ref{fig:erd-patient} & 患者主档与风险资料 & 6 & 患者、过敏、同意及复核事件\\
图~\ref{fig:erd-workforce} & 牙医、资质与治疗目录 & 6 & 任职、证书、项目及资质要求\\
图~\ref{fig:erd-scheduling} & 排班、资源与预约 & 10 & 时段、牙椅、预约与状态历史\\
图~\ref{fig:erd-clinical} & 临床就诊与病历版本 & 8 & 就诊、实施项目、病历与附件\\
图~\ref{fig:erd-billing} & 结算、收款与成本 & 7 & 实施项目—明细—收款—退款\\
图~\ref{fig:erd-financial} & 期间核算与激励 & 4 & 费用归属、期间快照与关闭异常\\
图~\ref{fig:erd-expansion} & 扩张情景与候选诊所 & 4 & 情景假设、候选点与预测结果\\
\bottomrule\end{tabular}\end{table}

\noindent{}阅读时先沿核心业务链定位，再进入分域图核对字段与基数。下表提供跨图检索路径；外键可空性、删除动作与业务门槛见表~\ref{tab:domain-key-01}--\ref{tab:domain-key-09}。
\begin{table}[H]\centering\small
\caption{全景关系图的跨域核对路径}\label{tab:erd-cross-domain-paths}
\begin{tabular}{@{}L{27mm}L{72mm}L{50mm}@{}}\toprule
业务交接 & 关键关系表 / Key Tables & 核对重点\\\midrule
账户授权 & \path{user_role_assignments} $\to$ \path{app_users}, \path{roles}, \path{clinics} & 岗位、诊所范围与授权人留痕\\
本人绑定 & \path{patient_account_links} $\to$ \path{patients}, \path{app_users} & 外部登录身份不直接替代患者编号\\
安排预约 & \path{appointments} $\to$ \path{patients}, \path{dentists}, \path{clinic_chairs} & 时段与牙椅冲突需事务校验\\
到诊交接 & \path{clinical_visits} $\to$ \path{appointments} & 患者、诊所与预约信息保持一致\\
实施治疗 & \path{actual_treatment_items} $\to$ \path{clinical_visits}, \path{treatment_catalog} & 计划项目与实际项目分开保存\\
逐项开单 & \path{invoice_lines} $\to$ \path{invoices}, \path{actual_treatment_items} & 明细回溯到实际实施项目\\
退款追溯 & \path{payment_refunds} $\to$ \path{payments} $\to$ \path{invoices} & 原收款、累计可退额与期间归属\\
期间核定 & \path{incentive_pools} $\to$ \path{financial_period_results} $\to$ \path{financial_reporting_periods} & 激励依据锁定的核算快照\\
\bottomrule\end{tabular}\end{table}

% Native-size vector ERD plates. The source of all ten plates is the same
% audited 61-table schema-model.json; no image rasterisation is involved.
\clearpage
\figerdplate{00_panorama}{ADC Nexus 全景关系数据图}{fig:erd-panorama}
\figerdplate{01_organization_access}{组织、账户与授权关系数据图}{fig:erd-organization}
\figerdplate{02_authentication_notifications}{身份认证与通知关系数据图}{fig:erd-authentication}
\figerdplate{03_patient_master}{患者主档与风险资料关系数据图}{fig:erd-patient}
\figerdplate{04_workforce_catalog}{牙医、资质与治疗目录关系数据图}{fig:erd-workforce}
\figerdplate{05_scheduling_resources}{排班、资源与预约关系数据图}{fig:erd-scheduling}
\figerdplate{06_clinical_care}{临床就诊与病历版本关系数据图}{fig:erd-clinical}
\figerdplate{07_billing_cash}{结算、收款与成本关系数据图}{fig:erd-billing}
\figerdplate{08_financial_management}{期间核算与激励关系数据图}{fig:erd-financial}
\figerdplate{09_expansion_planning}{扩张情景与候选诊所关系数据图}{fig:erd-expansion}
\clearpage


\subsection{领域对象到关系表的转换}
领域模型向关系模型转换时，业务行为由领域对象与业务服务承担，边界类、控制类不直接生成业务表。值对象依据查询、约束、版本及审计需要确定是否独立存储。角色—权限、治疗项目—资质要求等多对多关系使用桥接表；一对一及一对零至一关系同时设置外键和唯一约束。表~\ref{tab:class-table-map} 汇总主链的映射结果，随后逐域列明物理字段、可空性、删除动作与业务含义，供图文对照；完整字段定义见附属数据字典。

\begin{table}[!htbp]\centering\small
\caption{领域对象与关系表的对应}\label{tab:class-table-map}
\begin{tabular}{@{}L{27mm}L{68mm}L{47mm}@{}}\toprule
领域对象 & 主要关系表 & 转换时保留的约束\\\midrule
患者及风险 & \makecell[l]{\texttt{patients}\\\texttt{patient\_allergies}} & 内部主键与全网患者编号分离；过敏为一对多记录\\
预约及变更 & \makecell[l]{\texttt{appointments}\\\texttt{appointment\_status\_history}} & 预约编号唯一；当前状态与变更历史分开\\
就诊及治疗 & \makecell[l]{\texttt{clinical\_visits}\\\texttt{actual\_treatment\_items}} & 预约与就诊至多一对一；实施项目为一对多\\
结算及明细 & \makecell[l]{\texttt{invoices}\\\texttt{invoice\_lines}} & 就诊与单据至多一对一；治疗行不重复计费\\
收退款 & \makecell[l]{\texttt{payments}\\\texttt{payment\_refunds}} & 退款指向原收款；资金事实保留独立流水\\
期间及结果 & \makecell[l]{\texttt{financial\_reporting\_periods}\\\texttt{financial\_period\_results}} & 期间编号唯一；期间与范围共同限定结果\\
\bottomrule\end{tabular}\end{table}

\noindent\textbf{关系图阅读方法。}\quad 图中的表卡显示物理表和字段；以下\emph{Key Path} 按“子表外键 $\rightarrow$ 父表主键”书写。
\emph{Nullable} 描述子表外键是否可空，\emph{ON DELETE} 是现有快照的删除动作，均不等同于业务操作权限。
\emph{UNIQUE}、状态互斥、金额上限与时段不重叠另由唯一约束、事务或应用规则执行。

\begin{table}[!htbp]\centering\small
\caption{关系图与关系说明的阅读层次}\label{tab:erd-reading-levels}
\begin{tabular}{@{}L{30mm}L{50mm}L{69mm}@{}}\toprule
层次 / View & 回答的问题 & 阅读范围\\\midrule
全景 / Panorama & 有哪些业务簇和跨域依赖？ & 61 张表及跨域连线；用于总览，不在 A4 上逐字段阅读。\\
分域 / Domain ERD & 本域的表怎样相连？ & 每张分域图显示字段、PK/FK 及关系箭头，九域逐张放大检查。\\
关系说明 / Key Path & 外键方向和可空性是什么？ & 下列九张关系表列出现有约束及其业务语义；未选的完整关系见附属清单。\\
字段字典 / Data Dictionary & 每个字段的类型和默认值是什么？ & 独立附录逐字段列出类型、键、可空、默认值和索引。\\
\bottomrule\end{tabular}\end{table}

\subsubsection{组织、账户与授权：关键关系与约束}
组织边界与岗位授权分开保存。组织、诊所是权限范围的定位点，岗位角色通过中间表授予用户，操作人不能由患者资料代替。

\begingroup\footnotesize\setlength{\tabcolsep}{3pt}
\begin{longtable}{@{}L{43mm}L{32mm}L{24mm}L{49mm}@{}}
\caption{组织、账户与授权关系路径（Key Paths）}\label{tab:domain-key-01}\\\toprule
Child table / FK & Parent table / PK & Null; Delete & Business meaning / 业务含义\\\midrule\endfirsthead
\toprule Child table / FK & Parent table / PK & Null; Delete & Business meaning / 业务含义\\\midrule\endhead
\texttt{clinics}\newline\texttt{organization\_id} & \texttt{organizations}\newline\texttt{id} & NO\newline\texttt{RESTRICT} & 诊所归属组织，跨诊所授权须先确定组织范围。\\
\rowcolor{Pale}\texttt{role\_permissions}\newline\texttt{role\_id} & \texttt{roles}\newline\texttt{id} & NO\newline\texttt{CASCADE} & 角色与权限采用关联表，避免在角色记录中存放权限列表。\\
\texttt{role\_permissions}\newline\texttt{permission\_id} & \texttt{permissions}\newline\texttt{id} & NO\newline\texttt{CASCADE} & 同一权限可由多个角色复用。\\
\rowcolor{Pale}\texttt{user\_role\_assignments}\newline\texttt{user\_id} & \texttt{app\_users}\newline\texttt{id} & NO\newline\texttt{CASCADE} & 一个账户可持有多个角色授权，授权本身可独立留痕。\\
\texttt{user\_role\_assignments}\newline\texttt{clinic\_id} & \texttt{clinics}\newline\texttt{id} & YES\newline\texttt{RESTRICT} & 诊所范围可选；不应把空值自动解释为全网权限。\\
\rowcolor{Pale}\texttt{audit\_events}\newline\texttt{actor\_user\_id} & \texttt{app\_users}\newline\texttt{id} & YES\newline\texttt{SET NULL} & 审计事实与操作账户关联；账户停用不能抹去历史事件。\\
\bottomrule\end{longtable}\endgroup


\subsubsection{身份认证与通知：关键关系与约束}
身份凭证、会话和患者本人绑定采用不同生命周期；通知记录与业务预约关联，但通知成功不等于预约已确认。

\begingroup\footnotesize\setlength{\tabcolsep}{3pt}
\begin{longtable}{@{}L{43mm}L{32mm}L{24mm}L{49mm}@{}}
\caption{身份认证与通知关系路径（Key Paths）}\label{tab:domain-key-02}\\\toprule
Child table / FK & Parent table / PK & Null; Delete & Business meaning / 业务含义\\\midrule\endfirsthead
\toprule Child table / FK & Parent table / PK & Null; Delete & Business meaning / 业务含义\\\midrule\endhead
\texttt{auth\_identities}\newline\texttt{user\_id} & \texttt{app\_users}\newline\texttt{id} & NO\newline\texttt{CASCADE} & 第三方身份映射归属内部账户，外部标识不能充当业务主键。\\
\rowcolor{Pale}\texttt{auth\_sessions}\newline\texttt{user\_id} & \texttt{app\_users}\newline\texttt{id} & NO\newline\texttt{CASCADE} & 登录会话可以撤销，不能替代账户或患者主档。\\
\texttt{patient\_account\_links}\newline\texttt{patient\_id} & \texttt{patients}\newline\texttt{id} & NO\newline\texttt{CASCADE} & 患者主档与本人账户的绑定须经过身份核实。\\
\rowcolor{Pale}\texttt{patient\_account\_links}\newline\texttt{user\_id} & \texttt{app\_users}\newline\texttt{id} & NO\newline\texttt{CASCADE} & 账户与患者关系独立记录，便于处理解绑及复核。\\
\texttt{notification\_events}\newline\texttt{appointment\_id} & \texttt{appointments}\newline\texttt{id} & YES\newline\texttt{CASCADE} & 预约通知保留投递事实，发送失败不改变预约状态。\\
\bottomrule\end{longtable}\endgroup


\subsubsection{患者主档与风险资料：关键关系与约束}
患者全网主档是临床、预约与收费的共同锚点。紧急联系人、过敏和同意材料是可重复、可追溯的独立事实。

\begingroup\footnotesize\setlength{\tabcolsep}{3pt}
\begin{longtable}{@{}L{43mm}L{32mm}L{24mm}L{49mm}@{}}
\caption{患者主档与风险资料关系路径（Key Paths）}\label{tab:domain-key-03}\\\toprule
Child table / FK & Parent table / PK & Null; Delete & Business meaning / 业务含义\\\midrule\endfirsthead
\toprule Child table / FK & Parent table / PK & Null; Delete & Business meaning / 业务含义\\\midrule\endhead
\texttt{patients}\newline\texttt{primary\_clinic\_id} & \texttt{clinics}\newline\texttt{id} & YES\newline\texttt{SET NULL} & 首诊诊所为主档来源，不限制患者后续跨店就诊。\\
\rowcolor{Pale}\texttt{patient\_allergies}\newline\texttt{patient\_id} & \texttt{patients}\newline\texttt{id} & NO\newline\texttt{CASCADE} & 同一患者可记录多种过敏，复核记录不可覆盖原始事实。\\
\texttt{patient\_consents}\newline\texttt{patient\_id} & \texttt{patients}\newline\texttt{id} & NO\newline\texttt{CASCADE} & 每次同意材料与患者关联，并保留签署和适用范围。\\
\rowcolor{Pale}\texttt{emergency\_contacts}\newline\texttt{patient\_id} & \texttt{patients}\newline\texttt{id} & NO\newline\texttt{CASCADE} & 紧急联系人可重复，不能压成患者表的固定单字段。\\
\texttt{patient\_duplicate}\newline\texttt{reviews}\newline\texttt{source\_patient\_id} & \texttt{patients}\newline\texttt{id} & NO\newline\texttt{RESTRICT} & 疑似重复患者先进入复核流程，不自动合并主档。\\
\rowcolor{Pale}\texttt{patient\_timeline}\newline\texttt{events}\newline\texttt{clinic\_id} & \texttt{clinics}\newline\texttt{id} & YES\newline\texttt{SET NULL} & 时间线事件保留发生诊所，支持跨店追溯。\\
\bottomrule\end{longtable}\endgroup


\subsubsection{牙医、资质与治疗目录：关键关系与约束}
牙医任职、资质证书、治疗目录和项目所需资质分别建模；能否执业须同时考虑门店、项目和服务日期。

\begingroup\footnotesize\setlength{\tabcolsep}{3pt}
\begin{longtable}{@{}L{43mm}L{32mm}L{24mm}L{49mm}@{}}
\caption{牙医、资质与治疗目录关系路径（Key Paths）}\label{tab:domain-key-04}\\\toprule
Child table / FK & Parent table / PK & Null; Delete & Business meaning / 业务含义\\\midrule\endfirsthead
\toprule Child table / FK & Parent table / PK & Null; Delete & Business meaning / 业务含义\\\midrule\endhead
\texttt{dentist\_clinic}\newline\texttt{assignments}\newline\texttt{dentist\_id} & \texttt{dentists}\newline\texttt{id} & NO\newline\texttt{RESTRICT} & 任职关系独立于牙医主档，可表示跨店任职和有效期。\\
\rowcolor{Pale}\texttt{dentist\_clinic}\newline\texttt{assignments}\newline\texttt{clinic\_id} & \texttt{clinics}\newline\texttt{id} & NO\newline\texttt{RESTRICT} & 门店排班须核对该诊所任职关系。\\
\texttt{dentist}\newline\texttt{qualifications}\newline\texttt{qualification\_type\_id} & \texttt{qualification}\newline\texttt{types}\newline\texttt{id} & NO\newline\texttt{RESTRICT} & 证书类别与证书实例分开，便于到期和更新。\\
\rowcolor{Pale}\texttt{treatment}\newline\texttt{qualification}\newline\texttt{requirements}\newline\texttt{treatment\_id} & \texttt{treatment}\newline\texttt{catalog}\newline\texttt{id} & NO\newline\texttt{CASCADE} & 治疗项目与所需资质是多对多规则，不嵌入目录文本。\\
\texttt{treatment}\newline\texttt{qualification}\newline\texttt{requirements}\newline\texttt{qualification\_type\_id} & \texttt{qualification}\newline\texttt{types}\newline\texttt{id} & NO\newline\texttt{RESTRICT} & 同类资质要求可服务多个治疗项目。\\
\bottomrule\end{longtable}\endgroup


\subsubsection{排班、资源与预约：关键关系与约束}
可用时段由营业时间、停诊、牙医排班和牙椅占用共同决定；预约表只保存已确认的选择，冲突校验仍需事务执行。

\begingroup\footnotesize\setlength{\tabcolsep}{3pt}
\begin{longtable}{@{}L{43mm}L{32mm}L{24mm}L{49mm}@{}}
\caption{排班、资源与预约关系路径（Key Paths）}\label{tab:domain-key-05}\\\toprule
Child table / FK & Parent table / PK & Null; Delete & Business meaning / 业务含义\\\midrule\endfirsthead
\toprule Child table / FK & Parent table / PK & Null; Delete & Business meaning / 业务含义\\\midrule\endhead
\texttt{clinic\_chairs}\newline\texttt{clinic\_id} & \texttt{clinics}\newline\texttt{id} & NO\newline\texttt{RESTRICT} & 牙椅是诊所资源，分配前须核对所属门店。\\
\rowcolor{Pale}\texttt{dentist\_availability}\newline\texttt{rules}\newline\texttt{dentist\_id} & \texttt{dentists}\newline\texttt{id} & NO\newline\texttt{CASCADE} & 常规可用规则与单次预约分离。\\
\texttt{dentist\_schedule}\newline\texttt{exceptions}\newline\texttt{clinic\_id} & \texttt{clinics}\newline\texttt{id} & YES\newline\texttt{CASCADE} & 门店级排班例外可关闭某日资源。\\
\rowcolor{Pale}\texttt{appointments}\newline\texttt{patient\_id} & \texttt{patients}\newline\texttt{id} & NO\newline\texttt{RESTRICT} & 预约始终指向已识别患者，不能以姓名字符串代替关联。\\
\texttt{appointments}\newline\texttt{chair\_id} & \texttt{clinic\_chairs}\newline\texttt{id} & YES\newline\texttt{RESTRICT} & 牙椅为可选分配；空值不代表资源校验通过。\\
\rowcolor{Pale}\texttt{appointment\_status}\newline\texttt{history}\newline\texttt{appointment\_id} & \texttt{appointments}\newline\texttt{id} & NO\newline\texttt{CASCADE} & 状态变更追加留痕，不以当前状态覆盖历史。\\
\texttt{waitlist\_offers}\newline\texttt{waitlist\_entry\_id} & \texttt{waitlist\_entries}\newline\texttt{id} & NO\newline\texttt{CASCADE} & 候补邀约保留来源候补记录，便于追踪补位结果。\\
\bottomrule\end{longtable}\endgroup


\subsubsection{临床就诊与病历版本：关键关系与约束}
临床事实沿预约—就诊—实施项目展开。病历版本、过敏确认和附件独立保存，以区分草稿、签署和后续更正。

\begingroup\footnotesize\setlength{\tabcolsep}{3pt}
\begin{longtable}{@{}L{43mm}L{32mm}L{24mm}L{49mm}@{}}
\caption{临床就诊与病历版本关系路径（Key Paths）}\label{tab:domain-key-06}\\\toprule
Child table / FK & Parent table / PK & Null; Delete & Business meaning / 业务含义\\\midrule\endfirsthead
\toprule Child table / FK & Parent table / PK & Null; Delete & Business meaning / 业务含义\\\midrule\endhead
\texttt{clinical\_visits}\newline\texttt{appointment\_id} & \texttt{appointments}\newline\texttt{id} & NO\newline\texttt{RESTRICT} & 就诊回指预约，须另外使用唯一约束验证至多一次有效就诊。\\
\rowcolor{Pale}\texttt{clinical\_visits}\newline\texttt{patient\_id} & \texttt{patients}\newline\texttt{id} & NO\newline\texttt{RESTRICT} & 就诊患者应与预约患者一致，该跨表规则由事务核对。\\
\texttt{actual\_treatment}\newline\texttt{items}\newline\texttt{visit\_id} & \texttt{clinical\_visits}\newline\texttt{id} & NO\newline\texttt{CASCADE} & 一处就诊可产生多项实际治疗，不能只保留计划项目。\\
\rowcolor{Pale}\texttt{clinical\_note}\newline\texttt{versions}\newline\texttt{clinical\_note\_id} & \texttt{clinical\_notes}\newline\texttt{id} & NO\newline\texttt{CASCADE} & 病历每次修订形成版本，签署版本不原地覆盖。\\
\texttt{visit\_allergy}\newline\texttt{acknowledgements}\newline\texttt{allergy\_id} & \texttt{patient}\newline\texttt{allergies}\newline\texttt{id} & NO\newline\texttt{RESTRICT} & 到诊时对具体过敏记录作确认，不能只有笼统风险标志。\\
\bottomrule\end{longtable}\endgroup


\subsubsection{结算、收款与成本：关键关系与约束}
结算关系以实际治疗为来源。单据、明细、收款、退款和调整采用分表事实，净收款通过有效流水重算。

\begingroup\footnotesize\setlength{\tabcolsep}{3pt}
\begin{longtable}{@{}L{43mm}L{32mm}L{24mm}L{49mm}@{}}
\caption{结算、收款与成本关系路径（Key Paths）}\label{tab:domain-key-07}\\\toprule
Child table / FK & Parent table / PK & Null; Delete & Business meaning / 业务含义\\\midrule\endfirsthead
\toprule Child table / FK & Parent table / PK & Null; Delete & Business meaning / 业务含义\\\midrule\endhead
\texttt{invoices}\newline\texttt{visit\_id} & \texttt{clinical\_visits}\newline\texttt{id} & NO\newline\texttt{RESTRICT} & 一处就诊至多开具一张有效结算单，唯一性需落到约束。\\
\rowcolor{Pale}\texttt{invoice\_lines}\newline\texttt{actual\_treatment\_item}\newline\texttt{id} & \texttt{actual\_treatment}\newline\texttt{items}\newline\texttt{id} & NO\newline\texttt{RESTRICT} & 每条计费明细溯源实际治疗，防止仅凭目录价格开单。\\
\texttt{payments}\newline\texttt{invoice\_id} & \texttt{invoices}\newline\texttt{id} & NO\newline\texttt{RESTRICT} & 一张单据可有多笔收款，金额累计须核对余额。\\
\rowcolor{Pale}\texttt{payment\_refunds}\newline\texttt{payment\_id} & \texttt{payments}\newline\texttt{id} & NO\newline\texttt{RESTRICT} & 退款必须指向原收款，并校验累计可退余额。\\
\texttt{invoice\_adjustments}\newline\texttt{approved\_by\_user\_id} & \texttt{app\_users}\newline\texttt{id} & YES\newline\texttt{SET NULL} & 授权调整保留审批人；未批准不能进入正式核算。\\
\rowcolor{Pale}\texttt{treatment\_cost}\newline\texttt{entries}\newline\texttt{actual\_treatment\_item}\newline\texttt{id} & \texttt{actual\_treatment}\newline\texttt{items}\newline\texttt{id} & NO\newline\texttt{RESTRICT} & 用品成本与实施项目关联，不能由账单金额倒推。\\
\bottomrule\end{longtable}\endgroup


\subsubsection{期间核算与激励：关键关系与约束}
期间、门店和网络三个核算层级须分开。期间核定结果和激励池是快照，缺失成本或费用凭据时先生成关闭异常。

\begingroup\footnotesize\setlength{\tabcolsep}{3pt}
\begin{longtable}{@{}L{43mm}L{32mm}L{24mm}L{49mm}@{}}
\caption{期间核算与激励关系路径（Key Paths）}\label{tab:domain-key-08}\\\toprule
Child table / FK & Parent table / PK & Null; Delete & Business meaning / 业务含义\\\midrule\endfirsthead
\toprule Child table / FK & Parent table / PK & Null; Delete & Business meaning / 业务含义\\\midrule\endhead
\texttt{operating\_expenses}\newline\texttt{reporting\_period\_id} & \texttt{financial}\newline\texttt{reporting}\newline\texttt{periods}\newline\texttt{id} & NO\newline\texttt{RESTRICT} & 费用归属明确核算期间，不能只按录入日统计。\\
\rowcolor{Pale}\texttt{financial\_close}\newline\texttt{exceptions}\newline\texttt{reporting\_period\_id} & \texttt{financial}\newline\texttt{reporting}\newline\texttt{periods}\newline\texttt{id} & NO\newline\texttt{RESTRICT} & 关闭异常指向待核定期间，异常未消除不得伪造完整结果。\\
\texttt{financial\_period}\newline\texttt{results}\newline\texttt{clinic\_id} & \texttt{clinics}\newline\texttt{id} & YES\newline\texttt{RESTRICT} & 门店结果与网络结果的范围必须明确区分。\\
\rowcolor{Pale}\texttt{incentive\_pools}\newline\texttt{network\_result\_id} & \texttt{financial\_period}\newline\texttt{results}\newline\texttt{id} & NO\newline\texttt{RESTRICT} & 激励池依据已核定网络结果，而非实时估算。\\
\texttt{incentive\_pools}\newline\texttt{approved\_by\_user\_id} & \texttt{app\_users}\newline\texttt{id} & YES\newline\texttt{SET NULL} & 批准责任独立留存，计算结果不自动等于可发放结果。\\
\bottomrule\end{longtable}\endgroup


\subsubsection{扩张情景与候选诊所：关键关系与约束}
扩张测算是情景分析，不是已批准开店事实。假设、候选诊所和模拟结果分别保存，使每次预测可复算。

\begingroup\footnotesize\setlength{\tabcolsep}{3pt}
\begin{longtable}{@{}L{43mm}L{32mm}L{24mm}L{49mm}@{}}
\caption{扩张情景与候选诊所关系路径（Key Paths）}\label{tab:domain-key-09}\\\toprule
Child table / FK & Parent table / PK & Null; Delete & Business meaning / 业务含义\\\midrule\endfirsthead
\toprule Child table / FK & Parent table / PK & Null; Delete & Business meaning / 业务含义\\\midrule\endhead
\texttt{expansion\_scenarios}\newline\texttt{created\_by\_user\_id} & \texttt{app\_users}\newline\texttt{id} & YES\newline\texttt{SET NULL} & 情景创建人用于方案责任追溯。\\
\rowcolor{Pale}\texttt{expansion\_scenarios}\newline\texttt{approved\_by\_user\_id} & \texttt{app\_users}\newline\texttt{id} & YES\newline\texttt{SET NULL} & 审批人可选，未批准的情景不得标成已定方案。\\
\texttt{expansion\_assumptions}\newline\texttt{scenario\_id} & \texttt{expansion}\newline\texttt{scenarios}\newline\texttt{id} & NO\newline\texttt{CASCADE} & 测算参数归属特定情景，不覆盖别的版本。\\
\rowcolor{Pale}\texttt{expansion\_candidate}\newline\texttt{clinics}\newline\texttt{scenario\_id} & \texttt{expansion}\newline\texttt{scenarios}\newline\texttt{id} & NO\newline\texttt{CASCADE} & 一个情景可比较多个候选地点。\\
\texttt{expansion\_scenario}\newline\texttt{results}\newline\texttt{scenario\_id} & \texttt{expansion}\newline\texttt{scenarios}\newline\texttt{id} & NO\newline\texttt{CASCADE} & 预测结果与输入情景关联，展示时注明非实际经营数据。\\
\bottomrule\end{longtable}\endgroup

\noindent\emph{说明：}上表精选的是能够解释九张分域图主线的关系，不以节省篇幅隐藏其余外键。
附属关系清单记录现有 119 条外键的完整目标、字段和删除动作；
图中未显示的三个目标增强外键仍按本节后文的迁移建议处理，未混入现有关系。


\begingroup\scriptsize\setlength{\tabcolsep}{2.5pt}
\begin{longtable}{@{}C{8mm}L{27mm}L{42mm}C{10mm}C{10mm}L{49mm}@{}}
\caption{关系数据表目录（61 表）}\label{tab:relational-catalog}\\
\toprule 序 & 业务簇 & 物理表 & 字段 & 外键 & 业务责任\\\midrule\endfirsthead
\toprule 序 & 业务簇 & 物理表 & 字段 & 外键 & 业务责任\\\midrule\endhead
1 & 组织、账户与授权 & \path{account_invitations} & 10 & 2 & 受控创建账户或授予初始角色时的邀请及失效记录。 \\
2 & 组织、账户与授权 & \path{app_users} & 18 & 0 & 员工或患者入口所使用的系统账户，保存受保护的认证状态与界面偏好。 \\
3 & 组织、账户与授权 & \path{audit_events} & 15 & 2 & 关键读取、修改、拒绝和审批事件的追加式审计轨迹。 \\
4 & 组织、账户与授权 & \path{clinics} & 16 & 1 & 组织下的门店主数据，保存中英文名称、地址、联系方式和牙椅容量。 \\
5 & 组织、账户与授权 & \path{organizations} & 9 & 0 & 连锁运营主体的根主数据，统一默认时区、币种及状态。 \\
6 & 组织、账户与授权 & \path{permissions} & 6 & 0 & 按业务域定义的最小授权动作目录。 \\
7 & 组织、账户与授权 & \path{role_permissions} & 3 & 2 & 角色与权限的多对多桥接表，复合主键防止重复授权。 \\
8 & 组织、账户与授权 & \path{roles} & 10 & 0 & 可分配给账户的岗位或系统角色目录。 \\
9 & 组织、账户与授权 & \path{system_settings} & 5 & 2 & 组织范围的键值配置，记录最后修改人。 \\
10 & 组织、账户与授权 & \path{user_role_assignments} & 9 & 4 & 账户、角色与可选诊所范围的有效期授权记录。 \\
11 & 身份认证与通知 & \path{auth_identities} & 8 & 1 & 微信等外部身份提供方主体与本地账户的唯一映射。 \\
12 & 身份认证与通知 & \path{auth_oauth_states} & 7 & 0 & OAuth 发起与回调之间的一次性状态校验记录。 \\
13 & 身份认证与通知 & \path{auth_otp_challenges} & 11 & 0 & 短信验证码挑战的哈希、有效期、尝试次数和消费状态。 \\
14 & 身份认证与通知 & \path{auth_sessions} & 10 & 1 & 可撤销的登录会话和最小化终端指纹。 \\
15 & 身份认证与通知 & \path{notification_events} & 13 & 1 & 短信、微信或站内通知的请求、供应商回执和投递结果。 \\
16 & 身份认证与通知 & \path{patient_account_links} & 6 & 2 & 经核实的患者主档与本人门户账户的一对一绑定。 \\
17 & 患者主档与风险资料 & \path{emergency_contacts} & 10 & 1 & 患者可重复的紧急联系人及联系优先级。 \\
18 & 患者主档与风险资料 & \path{patient_allergies} & 14 & 2 & 过敏物、反应、严重度、状态和最近核实责任。 \\
19 & 患者主档与风险资料 & \path{patient_consents} & 9 & 1 & 患者对特定对象或用途作出的同意、撤回及证据。 \\
20 & 患者主档与风险资料 & \path{patient_duplicate_reviews} & 10 & 3 & 疑似重复患者主档的人工复核任务与处置结论。 \\
21 & 患者主档与风险资料 & \path{patient_timeline_events} & 15 & 2 & 跨预约、就诊和结算的患者事件索引，不复制病历全文。 \\
22 & 患者主档与风险资料 & \path{patients} & 16 & 1 & 患者全网主档；内部键、全网患者编号和联系方式相互分离。 \\
23 & 牙医、资质与治疗目录 & \path{dentist_clinic_assignments} & 9 & 2 & 牙医在诊所的任职关系、有效期及主要执业点。 \\
24 & 牙医、资质与治疗目录 & \path{dentist_qualifications} & 11 & 2 & 牙医持有的证书、编号、签发机构和有效期。 \\
25 & 牙医、资质与治疗目录 & \path{dentists} & 15 & 0 & 牙医员工主档及当前雇佣状态。 \\
26 & 牙医、资质与治疗目录 & \path{qualification_types} & 9 & 0 & 执业证书、专项能力等可复用的资质类型目录。 \\
27 & 牙医、资质与治疗目录 & \path{treatment_catalog} & 16 & 0 & 带版本和生效区间的治疗项目、时长、价格与估算成本目录。 \\
28 & 牙医、资质与治疗目录 & \path{treatment_qualification_requirements} & 4 & 2 & 治疗项目与所需资质类型之间的桥接关系。 \\
29 & 排班、资源与预约 & \path{appointment_status_history} & 10 & 2 & 预约每次状态变化、原因、操作者和版本的追加历史。 \\
30 & 排班、资源与预约 & \path{appointments} & 21 & 6 & 患者、诊所、牙医、牙椅、治疗项目及计划时段的预约主记录。 \\
31 & 排班、资源与预约 & \path{clinic_chairs} & 12 & 1 & 诊所内可预约牙椅、房间、设备和可访问性状态。 \\
32 & 排班、资源与预约 & \path{clinic_closures} & 9 & 1 & 节假日、维护或临时停诊等营业例外。 \\
33 & 排班、资源与预约 & \path{clinic_operating_hours} & 7 & 1 & 诊所按星期定义的常规营业时段。 \\
34 & 排班、资源与预约 & \path{dentist_availability_rules} & 11 & 2 & 牙医在特定诊所的周期性可用时间规则。 \\
35 & 排班、资源与预约 & \path{dentist_schedule_exceptions} & 10 & 2 & 请假、加班或临时调班等一次性排班例外。 \\
36 & 排班、资源与预约 & \path{requested_treatment_items} & 13 & 3 & 一次预约中患者提出的一个或多个治疗诉求。 \\
37 & 排班、资源与预约 & \path{waitlist_entries} & 13 & 3 & 患者对治疗、诊所和候选时段的候补需求。 \\
38 & 排班、资源与预约 & \path{waitlist_offers} & 14 & 5 & 空档向候补患者发出的限时邀约与接受结果。 \\
39 & 临床就诊与病历版本 & \path{actual_treatment_items} & 17 & 4 & 本次实际实施的目录项目、牙位、数量、操作者和时间。 \\
40 & 临床就诊与病历版本 & \path{clinical_attachments} & 13 & 1 & 与就诊关联的影像、文档或其他受控附件元数据。 \\
41 & 临床就诊与病历版本 & \path{clinical_note_versions} & 9 & 2 & 临床记录每次草拟、签署或修订的不可覆盖版本。 \\
42 & 临床就诊与病历版本 & \path{clinical_notes} & 9 & 2 & 一次就诊的临床记录容器及当前版本指针。 \\
43 & 临床就诊与病历版本 & \path{clinical_visits} & 17 & 6 & 由预约进入的一次临床就诊边界及风险复核状态。 \\
44 & 临床就诊与病历版本 & \path{odontogram_observations} & 15 & 3 & 按标准牙位和牙面记录的检查所见。 \\
45 & 临床就诊与病历版本 & \path{treatment_plans} & 14 & 1 & 就诊中形成的治疗方案版本、状态和临床说明。 \\
46 & 临床就诊与病历版本 & \path{visit_allergy_acknowledgements} & 9 & 3 & 牙医在本次就诊对具体过敏记录的复核与结论。 \\
47 & 结算、收款与成本 & \path{financial_reporting_periods} & 16 & 2 & 季度等财务边界、关闭清单、锁定和重开版本。 \\
48 & 结算、收款与成本 & \path{invoice_adjustments} & 10 & 3 & 经说明和审批的折让、补差或其他结算调整。 \\
49 & 结算、收款与成本 & \path{invoice_lines} & 15 & 2 & 来源于实际治疗项目的价格、数量和成本快照。 \\
50 & 结算、收款与成本 & \path{invoices} & 23 & 5 & 由已完成就诊形成的诊疗结算单及应收余额。 \\
51 & 结算、收款与成本 & \path{payment_refunds} & 10 & 2 & 指向原收款的退款流水，不覆盖原收款事实。 \\
52 & 结算、收款与成本 & \path{payments} & 11 & 2 & 已核实收款的独立流水及外部参考号。 \\
53 & 结算、收款与成本 & \path{treatment_cost_entries} & 13 & 4 & 按实际治疗项目归集的目录、材料、技工室和提供者成本。 \\
54 & 期间核算与激励 & \path{financial_close_exceptions} & 13 & 2 & 阻断期间关闭的数据缺项、差异或归属异常。 \\
55 & 期间核算与激励 & \path{financial_period_results} & 17 & 3 & 按期间和范围保存的核定输入快照与利润结果。 \\
56 & 期间核算与激励 & \path{incentive_pools} & 11 & 3 & 从已批准全网期间结果计算并审批的激励资金池。 \\
57 & 期间核算与激励 & \path{operating_expenses} & 13 & 4 & 诊所和财务期间范围内的运营费用及审批责任。 \\
58 & 扩张情景与候选诊所 & \path{expansion_assumptions} & 10 & 1 & 每项预测假设的值、单位和可追溯来源。 \\
59 & 扩张情景与候选诊所 & \path{expansion_candidate_clinics} & 20 & 1 & 情景中的候选地区、资源、需求与初始投资假设。 \\
60 & 扩张情景与候选诊所 & \path{expansion_scenario_results} & 15 & 1 & 扩张情景计算的收入、成本、利润、回收期与风险结果。 \\
61 & 扩张情景与候选诊所 & \path{expansion_scenarios} & 16 & 2 & 带版本、期限、免责声明和审批状态的扩张规划情景。 \\
\bottomrule\end{longtable}
\endgroup


\subsection{跨表完整性与结构演进}
数据库约束负责主键、外键、非空、唯一性及已声明的字段范围；资源不重叠、服务日资质有效、累计退款上限、病历完成门槛与期间关闭条件，则由事务和应用服务协同执行。表~\ref{tab:relational-rules} 列出 20 条跨表业务规则及其执行层，补充关系图未能单独表达的约束。

\begingroup\scriptsize\setlength{\tabcolsep}{3pt}
\begin{longtable}{@{}C{14mm}L{30mm}L{29mm}L{74mm}@{}}
\caption{跨表业务完整性规则}\label{tab:relational-rules}\\
\toprule 编号 & 规则 & 执行层 & 可核对的约束\\\midrule\endfirsthead
\toprule 编号 & 规则 & 执行层 & 可核对的约束\\\midrule\endhead
BR-01 & 全网患者编号唯一 & UNIQUE & patient\_code 在全网唯一且创建后不因电话、诊所或姓名变化而改写。 \\
BR-02 & 本人账户一对一绑定 & 双 UNIQUE + 核实流程 & 一个患者最多绑定一个本人门户账户，一个账户也最多绑定一个患者主档。 \\
BR-03 & 预约时段顺序 & 字段校验 + 服务 & scheduled\_end\_at 必须晚于 scheduled\_start\_at，时长与起止时间保持一致。 \\
BR-04 & 预约资源不重叠 & 事务内冲突查询 & 有效预约的同一牙医或同一牙椅在重叠时段只能被一条预约占用。SQLite/D1 需在提交事务中再次检查。 \\
BR-05 & 预约资格有效 & 应用服务 & 服务日必须处于牙医任职与所需资质的有效区间内。 \\
BR-06 & 营业与排班共同约束 & 候选查询 + 提交复核 & 候选时段必须同时满足门店营业、非停诊、牙医可用及无排班例外。 \\
BR-07 & 预约状态历史追加 & 同一事务 & 更新预约当前状态时必须追加对应历史，并保存原因、操作者和版本。 \\
BR-08 & 候补邀约并发控制 & 状态条件更新 + 事务 & 一个释放时段只能由一个仍有效的邀约成功接受，过期邀约不能创建预约。 \\
BR-09 & 预约与就诊至多一对一 & UNIQUE & 同一预约最多形成一次临床就诊；未到诊或取消预约可以没有就诊。 \\
BR-10 & 过敏风险按次复核 & 应用服务 & 开始治疗前需由本次牙医确认仍有效的过敏信息；历史确认不能替代本次复核。 \\
BR-11 & 计划与实际治疗分离 & 独立事实表 & 计划项目不自动等于已实施项目；结算只读取实际治疗事实。 \\
BR-12 & 临床记录版本不可覆盖 & 版本唯一 + 追加写入 & 已签署记录的修订创建新版本；旧版本、作者和签署时间保持可追溯。 \\
BR-13 & 一次就诊至多一张结算单 & UNIQUE & 重复生成请求返回既有结算单，不产生第二张有效单据。 \\
BR-14 & 实际治疗只计费一次 & UNIQUE & invoice\_lines.actual\_treatment\_item\_id 唯一，防止同一实施项目重复进入账单。 \\
BR-15 & 收款不超过应收余额 & 事务 + 金额校验 & 新增有效收款后累计净收款不得超过单据总额；并发收款必须串行核对余额。 \\
BR-16 & 退款不超过原收款 & 事务 + 金额校验 & 同一原收款的有效退款累计不得超过该收款金额，退款不删除或改写原收款。 \\
BR-17 & 期间关闭前无阻断异常 & 关闭服务 & 仍有 OPEN 且具阻断级别的异常时拒绝锁定期间。 \\
BR-18 & 期间范围结果唯一 & 复合 UNIQUE & 同一期间和 scope\_key 只有一份当前结果，核定时保存计算输入快照。 \\
BR-19 & 审计最小化 & 写入策略 & 审计记录操作主体、对象和结果，但不写明文密码、验证码、令牌或完整病历正文。 \\
BR-20 & 预测假设可追溯 & 版本 + 来源字段 & 扩张结果必须绑定具体情景版本和假设来源；批准前保持 DRAFT 并展示免责声明。 \\
\bottomrule\end{longtable}
\endgroup


当前结构另有三项外键增强建议：\path{account_invitations.role_id} 关联 \path{roles.id}，\path{account_invitations.clinic_id} 关联 \path{clinics.id}，\path{odontogram_observations.supersedes_observation_id} 建立到本表 \path{id} 的可空自引用。迁移前须检查孤儿值与删除策略，自引用还须排除循环链。这三项尚未实施，不计入 119 条现有外键。

结构校验还识别出 52 个外键列缺少前导索引的性能评估候选。是否增建索引，应结合 \texttt{EXPLAIN QUERY PLAN}、实际查询频率及写入成本确定。预约冲突、并发收退款与期间关闭另需业务验收，DDL 执行成功仅证明相应结构可建立。字段级字典、完整关系清单、候选索引脚本及可编辑 ERD 随 D 部分专项文件提供。

\subsection{查询与变更的处理}
前台搜索患者时先按授权诊所和必要字段缩小候选，再显示经遮罩的联系方式；跨店查找需要更高权限或明确的诊疗关联。预约日视图按诊所、日期和牙医索引查询，候选时段还要排除停诊、闭店及无效资质。临床详情以一次就诊为主线加载必要病史、观察和实施项目，避免在多个页面重复查询并使版本不一致。结算查询同时呈现应收、有效收款、退款和余额的计算依据。

更新依次执行版本读取、修改提交、当前版本核对与事件写入。过期版本提交时返回冲突，提示用户核对其间变更。季度关闭先形成异常清单，待费用及退款归属明确后再保存核定快照；已关闭期间以快照为查询依据，关闭后的新增事实由实时账本另行说明。

\subsection{季度核算口径}
牙医贡献需把已核实的收入、治疗用品成本与可归属运营费用关联到牙医及季度。设计口径为
\begin{equation}
P_{d,q}=R_{d,q}-M_{d,q}-O_{d,q},\qquad B_q=0.10\max(P_q,0),
\end{equation}
其中 $R_{d,q}$ 为牙医 $d$ 在季度 $q$ 的已确认收入，$M_{d,q}$ 为相应用品成本，$O_{d,q}$ 为按批准规则分配的运营费用，$P_q$ 为已核定的全网季度净利润。成本或费用缺少凭据时应标示为待归属，不把零值当作已核实。收入确认、退款跨期和共用费用分摊均须由机构确定制度后才能成为正式财务口径。现有程序尚未完成牙医级费用分摊，故不能宣称 $P_{d,q}$ 已完整可用。

\subsection{经营指标的定义}
预约、临床和财务指标应使用相同的日期、诊所与状态过滤范围。表~\ref{tab:metric-definitions} 约定报表展示时必须同时交代的分子、分母和排除项。它们是分析口径建议，不代表门店已经批准绩效考核规则。不同诊所的治疗类型、营业日和人员配置可能不同，横向比较时应同时显示这些背景条件。

\begin{longtable}{@{}L{30mm}L{47mm}L{48mm}L{23mm}@{}}
\caption{经营指标的计算与解释}\label{tab:metric-definitions}\\\toprule
指标 & 分子或取值 & 分母及范围 & 注意事项\\\midrule\endfirsthead
\toprule 指标 & 分子或取值 & 分母及范围 & 注意事项\\\midrule\endhead
主动取消率 & 在选定期间由患者或机构明确取消的预约数 & 同期原定应服务预约数；按机构规定的取消时点过滤 & 不含未到诊\\
未到诊率 & 到诊截止后确认未到诊的预约数 & 同期应到诊且未预先取消的预约数 & 到诊后未治疗另列\\
补位率 & 释放时段被候补确认并形成有效预约的次数 & 可用于候补的释放时段数 & 需保留原取消\\
牙椅利用率 & 牙椅被有效服务占用的分钟数 & 对应诊所实际营业且可用的牙椅分钟数 & 闭店日不入分母\\
临床完成率 & 满足完成条件的就诊次数 & 同期进入临床流程的有效就诊次数 & 草稿不得算完成\\
费用差异数 & 临床已实施项目与结算明细不一致的行数 & 同期已生成结算单的治疗行数 & 退单需留痕\\
资格覆盖率 & 指定治疗、日期和门店可排班的合格牙医数 & 对应治疗需求和时段；可按缺口数呈现 & 证书有效且可任职\\
季度净利润 & 核定收入减用品成本和运营费用 & 已关闭的组织与季度范围 & 不同于实时账本\\
\bottomrule\end{longtable}

分母为零时显示“不适用”。趋势图注明采用预约创建日、计划服务日或实际服务日，并按既定口径归属取消事件，防止跨期重复计入。财务数据同时标识期间是否关闭，当前估算与核定净利润分别呈现。改善幅度需在真实门店建立基线后测量。

\subsection{表示层：界面与操作组织}
界面按岗位任务而非数据库表组织：前台从患者档案进入预约和到诊；牙医从当日待诊队列进入病史复核、牙位观察与临床记录；收费员从已完成就诊进入逐项结算；经理查看本店运营，总部财务查看授权期间和核算异常。每个详情页展示编号、当前状态、最近修改时间及允许的下一步，减少跨页面核对。重要操作显示影响对象和确认信息；无权限、无数据与服务失败分别给出不同提示。

中英文切换保留当前页面、诊所范围、日期筛选和未提交内容。诊所、牙医和治疗目录的正式名称属于业务数据，不由语言切换任意翻译；字段标签、状态说明与错误提示则随界面语言变化。小屏幕优先保留关键任务，密集统计改为摘要和详情。状态始终配文字，颜色只是辅助，键盘焦点与表单错误应能被辨认。

\subsection{四类岗位的任务界面}
前台工作区默认显示当前诊所与日期，左侧安排搜索和日程，右侧呈现选中预约的患者、时间、牙医及牙椅。新预约先定位患者，确认页再展示资格与资源校验结果。取消或改期记录发起方及原因，完成后返回原日程并突出变动项，使前台在同一任务路径中完成核对。

牙医工作区按“待到诊—已到诊—进行中—待签署”分组，患者风险提示放在治疗录入之前。牙位图用于定位观察，旁边保留文本描述，防止只凭颜色或图形理解病情。计划、实际治疗和记录签署各有独立操作，页面应明确哪部分已经保存、哪部分只是草稿；离开页面前若有未提交修改须提醒。

收费工作区从已完成就诊进入，先列出实际治疗来源，再列出价格、数量和合计。收费员可以核对价格依据、登记收款或申请退款，但不能在此修改临床实施项目。患者询问费用时，系统应能从结算行返回相应治疗行和服务日期；若来源仍有疑问，退回临床核查而不是临时覆盖账单。

经理和财务工作区显示指标定义、日期、诊所范围与更新时间。取消率注明预约分母，门店损益说明待归属费用，已锁定季度结果与当前收退款账本分列。异常按阻断关闭、待人工复核、一般提示排序，并可进入对应记录，支持从汇总结果追查明细依据。

\begin{table}[!htbp]\centering\small
\caption{重点页面的反馈与防错设计}\label{tab:ui-feedback}
\begin{tabular}{@{}L{31mm}L{59mm}L{53mm}@{}}\toprule
任务 & 正常反馈 & 错误或例外反馈\\\midrule
预约提交 & 显示编号、诊所、时间、牙医及预约状态 & 列出冲突资源或无效资格；保留已填条件\\
临床完成 & 显示完成时间和签署版本 & 逐项指出过敏、同意或记录缺项\\
结算生成 & 显示来源就诊及逐项金额 & 重复开单返回已有编号；来源不全停止提交\\
退款登记 & 显示原收款与剩余可退额 & 超额或权限不足时不改变任何流水\\
季度关闭 & 显示输入版本和锁定人 & 缺项、版本变化或未核实费用阻断关闭\\
身份验证 & 成功后进入对应岗位首页 & 通道未配置、验证失败和账户停用分别处理\\
\bottomrule\end{tabular}\end{table}

\Needspace{7\baselineskip}
\section{实现核对与测试}\label{sec:implementation}
\subsection{现有系统与设计的对应}
现有程序使用 React、TypeScript 和 Cloudflare 兼容运行环境，数据层为 SQLite/D1。程序已提供患者查找与建档、重复候选复核、预约创建和状态变更、过敏确认、临床观察及记录、逐项结算与收退款、人员资质查询、经营分析、角色与诊所范围控制，以及中文和英文界面。系统处理的是课程用合成数据，不含真实患者资料。

目标方案与现有实现分别核对。候补邀约确认回写、临床计划及同意证据的完整维护、牙医级费用归属与季度贡献，以及正式短信、微信通道仍待补齐或联调。患者档案最终合并、病史完整版本管理和统一拒绝事件日志亦未列为已验收功能。现有签署记录操作账号、牙医、时间及版本，其效力限于系统内责任留痕；收退款由人工核实凭据后登记，程序不执行银行划转。

\subsection{敏感资料保护}
牙科病史、过敏信息、诊疗记录和联系方式不应因统一平台而被无限制复制。查询先确认访问目的及岗位，再限定诊所或治疗关系；患者端只显示经核实绑定的本人资料。列表默认展示必要摘要，手机号等联系信息按岗位和任务决定是否完整显示。导出、批量查询和账户授权比普通详情浏览风险更高，应单独记录操作者、范围与理由。

审计记录保存操作主体、对象、时间和处理结果，排除明文密码、一次性验证码、完整病历正文及未脱敏回调。账号停用、诊所调岗或密码重置后，应撤销已失效授权所对应的会话，并在接口层重新核验权限。开发与课程展示使用合成资料，公开测试环境不得导入真实患者记录。

正式上线前，机构须明确隐私告知、患者授权、资料保留期限、更正及删除请求的处理机制，并审查第三方服务合同和数据流向。这些事项构成实施的治理前提，尚需专门审查。跨店共享始终以诊疗必要性与有效授权为条件。

\subsection{部署前置条件}
部署环境需配置独立的应用域名、数据库、密钥与管理员初始账户。短信和微信通道必须填入机构批准的供应商凭据，验证回调域名、签名、限频及失败处理，再进行真实号码和微信账户的联调。没有这些外部条件时，密码登录仍可使用已建立的账号，但手机号验证码与微信登录不能被描述为已开通。

上线顺序宜先核对诊所、牙医、任职和治疗目录，再导入经清洗的患者资料；之后在少量门店试行预约、临床和结算，最后引入总部期间核算。每一步都要有回退条件：若患者身份映射无法核实，停止跨店合并；若治疗目录与旧收费项目不能对上，停止自动生成结算；若费用归属不完整，期间保持开放。正式切换前应由各岗位确认当日未完成事务和历史余额，以免新旧系统同时写入造成重复预约或重复收费。

\subsection{测试与验收重点}
测试以业务任务为单位，交叉核对界面、接口与持久化记录。拒绝路径重点确认原数据未被改写；并发、资源及资金约束采用专门场景验证。表~\ref{tab:tests} 同时标明预期结果和当前验收状态。

\begin{longtable}{@{}C{13mm}L{54mm}L{64mm}L{18mm}@{}}
\caption{关键业务验收场景}\label{tab:tests}\\\toprule
编号 & 操作与条件 & 预期结果 & 状态\\\midrule\endfirsthead
\toprule 编号 & 操作与条件 & 预期结果 & 状态\\\midrule\endhead
T-01 & 跨店检索相似患者并提交建档 & 提示可见范围内的重复线索，未经核实不合并 & 局部回归\\
T-02 & 无有效资质或牙椅冲突时预约 & 明确拒绝并保留原预约资料 & 待专项验收\\
T-03 & 同一资源并发提交两次 & 至多一笔预约成立 & 待专项验收\\
T-04 & 临床风险未复核或记录未签署 & 拒绝完成就诊，列出阻断项 & 局部回归\\
T-05 & 同一就诊重复生成结算单 & 返回既有单据，不重复计费 & 自动回归\\
T-06 & 超额收款或超额退款 & 拒绝并保持流水和余额不变 & 待专项验收\\
T-07 & 门店账号读取其他店明细 & 服务端拒绝，响应不泄露数据 & 自动回归\\
T-08 & 候补邀约被拒绝或逾期 & 不生成正式预约，保留处理记录 & 待开发验收\\
T-09 & 费用缺项时关闭季度 & 不锁定期间，列出缺失归属 & 待专项验收\\
T-10 & 未配置短信或微信供应商 & 不发放模拟验证码，不建立会话 & 配置核对\\
T-11 & 正式手机号找回密码 & 验证真实挑战，仅更改本人密码并撤销旧会话 & 待外部联调\\
T-12 & 中英文切换后继续当前任务 & 页面与筛选不丢失，权限不变 & 界面核对\\
\bottomrule\end{longtable}

2026 年 10 月 8 日的本地发布核对使用临时合成数据，代码检查、构建及 30 项 HTTP/数据库断言通过。该证据覆盖选定操作与拒绝路径；生产并发、真实短信送达、微信回调及管理成效仍需分别验证。正式上线前应完成供应商联调、隐私告知、恢复演练、岗位试用和负载试验。性能评估采用接近 30 家诊所规模的数据量，记录第 95 百分位响应时间及相应负载条件。

\subsection{后续验证的取样办法}
预约并发试验使用两名前台账号同时请求相同牙医和牙椅、部分重叠的时段。除检查最终仅有一笔有效预约，还应检查失败请求是否留下状态历史、通知或候补占用。资质试验至少包含有效、当天到期、前一天过期、尚未生效四种日期。临床试验分别省略过敏复核、同意证据、实施项目和签署，以确认完成门槛给出不同、可执行的修正提示。

财务试验需构造部分收款、重复收款、全额退款和跨季度退款，比较单据余额、收退款流水与期间结果。季度关闭应覆盖费用缺项、核算过程中输入变化、关闭后普通修改、授权重开四类情况；核对的重点不是只看页面上的总额，还要追查每个差异来自哪张结算单或哪笔费用。患者端测试使用两个已绑定账户交换记录编号，服务端必须拒绝越权读取。

岗位试用分别安排前台查找旧患者并改期、牙医复核过敏并签署治疗记录、收费员解释账单并登记部分收款。记录完成时间、求助次数、错误回退及额外口头说明，与原工作方式按相同口径比较。重复建立线下台账的情况也应纳入观察，用于识别尚未消除的信息交接负担，并据此修订用例步骤和界面文案。

\section{结论}
本设计围绕跨店患者识别、临床治疗与收费衔接、门店账务与季度核算三处关键交接，建立了相互对应的岗位职责、对象关系和数据约束。业务用例界定服务结果，系统用例落实操作目标，序列图、类图与状态机共同说明职责分配及合法变化，关系模型则保存可复核的业务依据。现有课程程序已支撑主要登记、预约、诊疗和结算流程。后续实施应优先完善费用归属与正式身份通道，验证预约并发和恢复能力，再通过门店岗位试用检验其对日常工作及经营核算的实际价值。

\clearpage
\phantomsection\addcontentsline{toc}{section}{参考文献}
\begin{thebibliography}{9}\small
\bibitem{case} Chen Z. \emph{Andrew's Dental Chain: Data Modelling}. W44694. Ivey Publishing, 2024，课程案例集。
\bibitem{course}《系统分析与设计》课程课件，精简版，2026 年 9 月。
\bibitem{omg} Object Management Group. \emph{OMG Unified Modeling Language Version 2.5.1}. \url{https://www.omg.org/spec/UML/2.5.1/PDF}
\bibitem{plantuml-ie} PlantUML. \emph{Information Engineering Diagram}. \url{https://plantuml.com/ie-diagram}
\bibitem{ada-history} American Dental Association. \emph{Medical/Dental Health History}. \url{https://www.ada.org/resources/practice/practice-management/medical-dental-health-history}
\bibitem{ada-records} American Dental Association. \emph{Documentation/Patient Records}. \url{https://www.ada.org/resources/practice/practice-management/documentation-patient-records}
\end{thebibliography}
\end{document}
